\documentclass[a4paper,11pt]{book}
%\documentclass[a4paper,twoside,11pt,titlepage]{book}
\usepackage{listings}
\usepackage[utf8]{inputenc}
\usepackage[spanish, es-tabla]{babel}

% \usepackage[style=list, number=none]{glossary} %
%\usepackage{titlesec}
%\usepackage{pailatino}

\decimalpoint
\usepackage{dcolumn}
\newcolumntype{.}{D{.}{\esperiod}{-1}}
\makeatletter
\addto\shorthandsspanish{\let\esperiod\es@period@code}
\makeatother

%\usepackage[chapter]{algorithm}
\RequirePackage{verbatim}
%\RequirePackage[Glenn]{fncychap}
\usepackage{fancyhdr}
\usepackage{graphicx}
\usepackage{afterpage}

\usepackage{longtable}

\usepackage[colorlinks=true,linkcolor=black,
			citecolor=blue, urlcolor=cyan]{hyperref} %referencia

\usepackage{float}
\newcommand*{\at}{@}

% ********************************************************************
% Re-usable information
% ********************************************************************

\newcommand{\myTitle}{Evaluación del estado del arte en técnicas estadísticas para el análisis comparativo de algoritmos de aprendizaje automático}
\newcommand{\myDegree}{Grado en Ingeniería Informática y Matemáticas}
\newcommand{\myName}{Jacinto Carrasco Castillo}
\newcommand{\myProf}{Francisco Herrera Triguero}
\newcommand{\myOtherProf}{María del Mar Rueda García}
%\newcommand{\mySupervisor}{Put name here\xspace}
\newcommand{\myFaculty}{Escuela Técnica Superior de Ingenierías Informática y de Telecomunicación, Facultad de Ciencias}
\newcommand{\myFacultyShort}{}
\newcommand{\myDepartment}{}
\newcommand{\myUni}{\protect{Universidad de Granada}}
\newcommand{\myLocation}{Granada}
\newcommand{\myTime}{Septiembre, 2016}
\newcommand{\myVersion}{Version 0.1}


\hypersetup{
pdfauthor = {\myName  (jacintocc\at correougr.es)},
pdftitle = {\myTitle},
pdfsubject = {},
pdfkeywords = {TFG, test estadisticos, aprendizaje automatico},
pdfproducer = {pdflatex}
}

% ********************************************************************
%\hyphenation{}


%\usepackage{doxygen/doxygen}
%\usepackage{pdfpages}
\usepackage{url}
\usepackage{colortbl,longtable}
\usepackage[stable]{footmisc}
\usepackage{amsthm}
%\usepackage{index}

%\makeindex
%\usepackage[style=long, cols=2,border=plain,toc=true,number=none]{glossary}
% \makeglossary

% Definición de comandos que me son útiles:
%\renewcommand{\indexname}{Índice alfabético}
%\renewcommand{\glossaryname}{Glosario}

\pagestyle{fancy}
\fancyhf{}
\fancyhead[LO]{\leftmark}
\fancyhead[RE]{\rightmark}
\fancyhead[RO,LE]{\textbf{\thepage}}
\renewcommand{\chaptermark}[1]{\markboth{\textbf{#1}}{}}
\renewcommand{\sectionmark}[1]{\markright{\textbf{\thesection. #1}}}

\setlength{\headheight}{1.5\headheight}

\newcommand{\HRule}{\rule{\linewidth}{0.5mm}}
%Definimos los tipos teorema, ejemplo y definición podremos usar estos tipos
%simplemente poniendo \begin{teorema} \end{teorema} ...
\newtheorem{ejemplo}{Ejemplo}[chapter]
\theoremstyle{plain}
\newtheorem{teorema}{Teorema}[chapter]
\newtheorem{lema}[teorema]{Lema}
\newtheorem{proposicion}[teorema]{Proposición}
\newtheorem{corolario}[teorema]{Corolario}

% Definición de estilo de definición



\newtheoremstyle{definition_wo_parentheses}
  {\topsep}% measure of space to leave above the theorem. E.g.: 3pt
  {\topsep}% measure of space to leave below the theorem. E.g.: 3pt
  {}% name of font to use in the body of the theorem
  {0pt}% measure of space to indent
  {\bfseries}% name of head font
  {.}% punctuation between head and body
  { }% space after theorem head; " " = normal interword space
  {\thmname{#1}\thmnumber{ #2.}\thmnote{ #3}}
  
  
\theoremstyle{definition_wo_parentheses}	
\newtheorem{definicion}{Definición}[chapter]

\definecolor{gray97}{gray}{.97}
\definecolor{gray75}{gray}{.75}
\definecolor{gray45}{gray}{.45}
\definecolor{gray30}{gray}{.30}

\lstset{ frame=Ltb,
     framerule=0.5pt,
     aboveskip=0.5cm,
     framextopmargin=3pt,
     framexbottommargin=3pt,
     framexleftmargin=0.1cm,
     framesep=0pt,
     rulesep=.4pt,
     backgroundcolor=\color{gray97},
     rulesepcolor=\color{black},
     %
     stringstyle=\ttfamily,
     showstringspaces = false,
     basicstyle=\scriptsize\ttfamily,
     commentstyle=\color{gray45},
     keywordstyle=\bfseries,
     %
     numbers=left,
     numbersep=6pt,
     numberstyle=\tiny,
     numberfirstline = false,
     breaklines=true,
   }
 
% minimizar fragmentado de listados
\lstnewenvironment{listing}[1][]
   {\lstset{#1}\pagebreak[0]}{\pagebreak[0]}

\lstdefinestyle{CodigoC}
   {
	basicstyle=\scriptsize,
	frame=single,
	language=C,
	numbers=left
   }
\lstdefinestyle{CodigoC++}
   {
	basicstyle=\small,
	frame=single,
	backgroundcolor=\color{gray30},
	language=C++,
	numbers=left
   }
\lstdefinestyle{CodigoR}
   {
	basicstyle=\small,
	frame=single,
	backgroundcolor=\color{gray30},
	language=R,
	numbers=left
   }
 
\lstdefinestyle{Consola}
   {basicstyle=\scriptsize\bf\ttfamily,
    backgroundcolor=\color{gray30},
    frame=single,
    numbers=none
   }


\newcommand{\bigrule}{\titlerule[0.5mm]}


%Para conseguir que en las páginas en blanco no ponga cabecerass
\makeatletter
\def\clearpage{%
  \ifvmode
    \ifnum \@dbltopnum =\m@ne
      \ifdim \pagetotal <\topskip
        \hbox{}
      \fi
    \fi
  \fi
  \newpage
  \thispagestyle{empty}
  \write\m@ne{}
  \vbox{}
  \penalty -\@Mi
}
\makeatother

\usepackage{pdfpages}
\usepackage{textcomp}
\usepackage{gensymb}
\usepackage{amsfonts}
\usepackage{amsmath}
\usepackage{amssymb}
\usepackage{enumerate}
\usepackage{multirow}
\usepackage{booktabs}
\usepackage{tabularx}

% Comandos para la escritura 
\newcommand{\abs}[1]{\left\lvert#1\right\rvert}
\newcommand{\norm}[1]{\left\lVert#1\right\rVert}
\newcommand{\var}{\text{var}}
\newcommand{\cov}{\text{cov}}
\newcommand\fnurl[2]{%
\href{#2}{#1}\footnote{\url{#2}}%
}


% Pseudocódigo
\usepackage{algorithm}
\usepackage{algorithmic}

% ****************************************************************
% DOCUMENTO
% ****************************************************************

\begin{document}
\pagenumbering{roman}
\begin{titlepage}

\newlength{\centeroffset}
\setlength{\centeroffset}{-0.5\oddsidemargin}
\addtolength{\centeroffset}{0.5\evensidemargin}
\thispagestyle{empty}

\noindent\hspace*{\centeroffset}\begin{minipage}{\textwidth}

\centering
\includegraphics[width=0.9\textwidth]{imagenes/logo_ugr.jpg}\\[1.4cm]

\textsc{ \Large TRABAJO FIN DE GRADO\\[0.2cm]}
\textsc{ \myDegree }\\[1cm]
% Upper part of the page
% 
% Title
{\Large \bfseries{ \myTitle }\\ }
\noindent\rule[-1ex]{\textwidth}{3pt}\\[3.5ex]
{\large\bfseries}
\end{minipage}

\vspace{2.5cm}
\noindent\hspace*{\centeroffset}\begin{minipage}{\textwidth}
\centering

\textbf{Autor}\\ {\myName}\\[2.5ex]
\textbf{Directores}\\
{\myProf\\
 \myOtherProf}\\[2cm]
\includegraphics[width=0.3\textwidth]{imagenes/etsiit_logo.png}\\[0.1cm]
\textsc{ \myFaculty }\\
\textsc{---}\\
\myTime
\end{minipage}
%\addtolength{\textwidth}{\centeroffset}
%\vspace{\stretch{2}}
\end{titlepage}



\chapter*{}
%\thispagestyle{empty}
%\cleardoublepage

%\thispagestyle{empty}

\begin{titlepage}
 
 
\setlength{\centeroffset}{-0.5\oddsidemargin}
\addtolength{\centeroffset}{0.5\evensidemargin}
\thispagestyle{empty}

\noindent\hspace*{\centeroffset}\begin{minipage}{\textwidth}

\centering
%\includegraphics[width=0.9\textwidth]{imagenes/logo_ugr.jpg}\\[1.4cm]

%\textsc{ \Large PROYECTO FIN DE CARRERA\\[0.2cm]}
%\textsc{ INGENIERÍA EN INFORMÁTICA}\\[1cm]
% Upper part of the page
% 

 \vspace{3.3cm}

%si el proyecto tiene logo poner aquí
\includegraphics{imagenes/logo.png} 
 \vspace{0.5cm}

% Title

{\Huge\bfseries Título del proyecto\\
}
\noindent\rule[-1ex]{\textwidth}{3pt}\\[3.5ex]
{\large\bfseries Subtítulo del proyecto.\\[4cm]}
\end{minipage}

\vspace{2.5cm}
\noindent\hspace*{\centeroffset}\begin{minipage}{\textwidth}
\centering

\textbf{Autor}\\ {Nombre Apellido1 Apellido2 (alumno)}\\[2.5ex]
\textbf{Directores}\\
{Nombre Apellido1 Apellido2 (tutor1)\\
Nombre Apellido1 Apellido2 (tutor2)}\\[2cm]
%\includegraphics[width=0.15\textwidth]{imagenes/tstc.png}\\[0.1cm]
%\textsc{Departamento de Teoría de la Señal, Telemática y Comunicaciones}\\
%\textsc{---}\\
%Granada, mes de 201
\end{minipage}
%\addtolength{\textwidth}{\centeroffset}
\vspace{\stretch{2}}

 
\end{titlepage}






\cleardoublepage
\thispagestyle{empty}

\begin{center}
{\large\bfseries \myTitle }\\
\end{center}
\begin{center}	
	\myName \\
\end{center}

%\vspace{0.7cm}
\noindent{\textbf{Palabras clave}: palabra\_clave1, palabra\_clave2, palabra\_clave3, ......}\\

\vspace{0.7cm}
\noindent{\textbf{Resumen}}\\

Poner aquí el resumen.
\cleardoublepage


\thispagestyle{empty}


\begin{center}
{\large\bfseries Project Title: Project Subtitle}\\
\end{center}
\begin{center}
	\myName	\\
\end{center}

%\vspace{0.7cm}
\noindent{\textbf{Keywords}: Keyword1, Keyword2, Keyword3, ....}\\

\vspace{0.7cm}
\noindent{\textbf{Abstract}}\\

Write here the abstract in English.

\chapter*{}
\thispagestyle{empty}

\noindent\rule[-1ex]{\textwidth}{2pt}\\[4.5ex]

Yo, \textbf{\myName}, alumno de la titulación \myDegree de la \textbf{\myFaculty de la \myUni}, con DNI 32056356, autorizo la
ubicación de la siguiente copia de mi Trabajo Fin de Grado en la biblioteca de ambos centros para que pueda ser
consultada por las personas que lo deseen.

\vspace{6cm}

\noindent Fdo: \myName

\vspace{2cm}

\begin{flushright}
Granada a X de septiembre de 2016.
\end{flushright}


\chapter*{}
\thispagestyle{empty}

\noindent\rule[-1ex]{\textwidth}{2pt}\\[4.5ex]

D. \textbf{ \myProf }, Profesor del Departamento de Ciencias de la Computación e Inteligencia Artificial de la Universidad de Granada.

\vspace{0.5cm}

D. \textbf{ \myOtherProf}, Profesora del Departamento de Estadística e Investigación Operativa de la Universidad de Granada.


\vspace{0.5cm}

\textbf{Informan:}

\vspace{0.5cm}

Que el presente trabajo, titulado \textit{\textbf{\myTitle}},
ha sido realizado bajo su supervisión por \textbf{\myName}, y autorizamos la defensa de dicho trabajo ante el tribunal
que corresponda.

\vspace{0.5cm}

Y para que conste, expiden y firman el presente informe en Granada a X de septiembre de 2016.

\vspace{1cm}

\textbf{Los directores:}

\vspace{5cm}

\noindent \textbf{\myProf \ \ \ \ \ \myOtherProf}

\chapter*{Agradecimientos}
\thispagestyle{empty}

       \vspace{1cm}


Poner aquí agradecimientos...


%\frontmatter
\tableofcontents
%\listoffigures
%\listoftables
%
%\mainmatter
%\setlength{\parskip}{5pt}

\chapter*{Planteamiento}\addcontentsline{toc}{chapter}{Planteamiento}
\markboth{Planteamiento}{}
\pagenumbering{arabic}
\setcounter{page}{1}

\section*{Contextualización}\addcontentsline{toc}{section}{Contextualización}
	
	El interés creciente en el campo del aprendizaje
automático ha significado un aumento en la propuesta
de nuevos algoritmos que pretenden mejorar a los 
ya existentes. Se describe a continuación el proceso
general para poder comparar algoritmos, basándonos en
las recomendaciones realizadas por Alpaydin 
(\cite{Alpaydin:2010:IML:1734076}) y Japkowicz  
(\cite{DBLP:books/cu/Japkowicz2011}).\\
	A la hora de evaluar un algoritmo,
nos interesa conocer el error asociado a un problema y 
su relación con respecto a otros. Sin embargo, debemos ser 
conscientes de que para ello debemos diseñar y realizar un 
experimento cuyos resultados sean fiables y minimizar el 
efecto producido por las circunstancias en las que se 
realizan o los datos disponibles. Por ejemplo, el error 
obtenido en los datos utilizados para entrenar el modelo 
(error de entrenamiento) será menor que el error en unos 
datos que no han participado en este entrenamiento y son 
usados para evaluar el modelo (error de test). Cabe mencionar 
también que el error de entrenamiento no es comparable debido 
a que modelos más complejos se ajustarán más a los datos de 
entrenamiento, sin que eso signifique que el modelo se 
comporte mejor de manera global. Hay algoritmos que se 
comportan de manera no determinista, y por tanto deseemos 
realizar varias ejecuciones para estimar circunstancias 
aleatorias. En definitiva, debemos basar nuestra evaluación 
del algoritmo en la distribución (desconocida) del error. 
Para una correcta estimación del error,  es importante tener 
en cuenta lo siguiente:
	
	\begin{itemize}
	\item El error depende del problema. El teorema de 
		\textit{No Free Lunch} (Wolpert y Macready, 
		\cite{Wolpert:1997:NFL:2221336.2221408}) 
		afirma que dados dos 
		algoritmos su rendimiento medio para todos los 
		posibles problemas es el mismo, lo que se traduce en 
		que según el problema se puede dar que un algoritmo 
		sea mejor que otro o al contrario.
	\item La partición de la base de datos referente al 
		problema en datos de entrenamiento y de test se 
		realiza para obtener una idea más acertada del error 
		cometido por el algoritmo para el problema. Sin 
		embargo, una vez que se considere válido el algoritmo 
		se entrenará con todos los datos disponibles.
	\item Para validar el algoritmo, se deben utilizar datos 
		que no hayan sido utilizados anteriormente de ninguna 
		manera.
	\item Normalmente, para la comparación del rendimiento de 
		algoritmos se usa el error cometido en la 
		clasificación, aunque se podrían considerar otros 
		factores como el tiempo necesario para entrenar el 
		algoritmo, su simplicidad...
	\end{itemize}
	
\subsection*{Diseño de experimentos}

	Al hablar de experimentos nos referimos al test o a la 
serie de test donde se manejan factores que afectan a la 
salida. Pretendemos minimizar aquellos factores que no se 
controlan y obtener resultados estadísticamente 
significativos. Para ajustar parámetros propios del modelo de 
aprendizaje automático que estemos utilizando existen 
diferentes aproximaciones: la \textit{mejor suposición} 
(donde quien realiza el experimento ajusta los parámetros 
basándose en su conocimiento sobre el problema), \textit{un 
factor cada vez} (se supone que los factores son 
independientes y se prueban diferentes valores para cada 
parámetro de uno en uno, mientras los demás se mantienen en 
una posición base) y \textit{diseño factorial} (se crea una 
rejilla con diferentes posibles valores para cada parámetro y 
se prueba en ellos, es más costoso).\\

	Para minimizar el número de ejecuciones necesarias, se 
suele utilizar conocimiento previo para reducir las 
combinaciones a probar. Otra estrategia consiste en el diseño 
de la superficie de respuesta. Se considera la siguiente 
situación:
		\[ r = g( f_1, \dots, f_F | \phi ), \]
	donde $r$ es la respuesta, $g$ el modelo utilizado, 
$f_1, \dots f_F$ los factores y $\phi$ la estimación empírica 
del modelo para cada configuración particular probada. El 
procedimiento consiste en ir añadiendo a $\phi$ las 
evaluaciones realizadas, ajustar $g$ según $\phi$ y buscar su 
máximo (o mínimo), evaluar el algoritmo en ese punto y 
ajustar $\phi$.\\
	En la realización de experimentos es necesaria la 
aleatorización a la hora de probar las diferentes 
configuraciones de parámetros (no tanto en nuestro contexto, 
pero el orden de las configuraciones al realizar experimentos 
donde interviene maquinaria por ejemplo podría interferir en los 
resultados debido a factores como la temperatura de la 
máquina). Otro factor a tener en cuenta es que es necesaria 
la repetición de las ejecuciones y promediar los resultados 
para así reducir el impacto de factores aleatorios. Para 
reducir la variabilidad de factores que no dependen del 
propio algoritmo a evaluar, como las bases de datos o las 
particiones realizadas, en la comparación de varios 
algoritmos utilizamos las mismas particiones para cada uno.

\paragraph{Directrices a la hora de realizar un experimento}
	\begin{enumerate}
	\item Fijar la intención del estudio
	\item Seleccionar las variables de respuesta
	\item Seleccionar los factores con los que se realizará 
		el experimento y los niveles a comprobar
	\item Diseño del experimento: La partición para 
		entrenamiento y test depende del tamaño de la base de 
		datos. Si la base de datos es pequeña puede haber una 
		alta variabilidad y obtener resultados no 
		concluyentes.
	\item Efectuar el experimento: Antes de llevar a cabo un 
		experimento de gran magnitud es aconsejable probar 
		pequeños conjuntos de datos para probar que funciona. 
		También es interesante guardar resultados intermedios 
		o las semillas aleatorias para facilitar la 
		reproducción de los datos.
	\item El experimentador debe ser imparcial y juzgar de 
		igual manera un algoritmo que otro, realizar 
		documentación ...
	\item Análisis estadístico: las afirmaciones y preguntas 
		que se realicen debe sostenerse estadísticamente 
		hablando. Es aconsejable un análisis visual para 
		exponer los datos obtenidos.
	\end{enumerate}
	
\subsection*{Test de hipótesis}

	Una vez se han extraído los datos, pretendemos obtener
las conclusiones de manera estadística. Para ello, la 
forma habitual de proceder es mediante los test de hipótesis.
En la realización de los test de hipótesis, el 
procedimiento es el siguiente: se define un estadístico que 
sigue una distribución conocida en el caso de que se cumpla 
la hipótesis realizada. Entonces si el estadístico calculado 
de la muestra tiene poca probabilidad de obtenerse de la 
distribución, se rechaza la hipótesis; en caso contrario, no. 
\\
	Surgen entonces las cuestiones: ¿Se pueden justificar los 
resultados estadísticamente o se obtienen por azar? ¿Son los 
conjuntos de datos representativos para el problema? Estas 
preguntas no se pueden responder de forma exhaustiva, es la 
labor de los test estadísticos reunir los datos disponibles 
para justificar la consistencia de las conclusiones. \\
	A la hora de realizar los test estadísticos podemos 
realizar las siguientes comparaciones:
	\begin{enumerate}
	\item Comparación de dos algoritmos en un dominio: $t$-
		test (paramétrico), test de McNemar (no paramétrico).
	\item Comparación de dos algoritmos en varios dominios: 
		\textit{Signed-Rank test} de Wilcoxon y test de signo 
		(no paramétricos).
	\item Comparación de múltiples algoritmos en varios 
		dominios: ANOVA (paramétrico), test de Friedman o 
		basados en permutaciones (no paramétricos). 
	\end{enumerate}

	Nótese que el test de \textit{Signed-Rank} de Wilcoxon 
para muestras apareadas puede aplicarse tanto para la 
comparación de dos como entre varios algoritmos en un 
dominio. Los test para varios algoritmos en varias bases de 
datos necesitan un test posterior en caso de rechazar la 
hipótesis de que todos los algoritmos se comporten de igual 
manera. Ejemplos de estos test son el de Tukey, de 
Bonferroni-Dunn o de Nemenyi.

\section*{Descripción del problema abordado}\addcontentsline{toc}{section}{Descripción del problema}

	El problema abordado consiste en el estudio de las 
técnicas estadísticas más relevantes en la comparación
del rendimiento de algoritmos de aprendizaje automático. 
En el estudio teórico se hace un recorrido por los test 
de hipótesis nula. En primer lugar, se describen y 
se incluyen los test paramétricos, ampliamente conocidos 
y utilizados en este contexto. Como alternativa a los 
test paramétricos se presentan los test no paramétricos, 
indicando las diferencias con los primeros, sus ventajas
e inconvenientes, las situaciones en las que son más 
favorables y un conjunto de test no paramétricos junto con su 
descripción y forma de usarlos. Una familia concreta de 
este tipo de test son los test basados en permutaciones,
de los cuales se presenta la base matemática asociada. Por
último, y como alternativa a los test de hipótesis nula, 
se presentan los test bayesianos, que utilizan la estadística 
bayesiana para la realización de test estadísticos para
la comparación de algoritmos. Al ser las técnicas
más novedosas, se presenta la revisión de varios
artículos recientes que aplican test no paramétricos y test
bayesianos al problema en cuestión.\\
	La vertiente práctica del problema consiste en la 
aplicación de estos test sobre los resultados obtenidos
de la evaluación de unos algoritmos en una serie de
conjuntos de datos. Este estudio práctico se ha llevado
a cabo en \texttt{R}, implementando un \textit{wrapper}
para el paquete de Java \texttt{JavaNPST}, que permite
la realización de test no paramétricos.

\section*{Técnicas utilizadas}\addcontentsline{toc}{section}{Técnicas utilizadas}

	Las principal área matemática presente en este trabajo es
la inferencia estadística, de la que se utilizan conceptos
y resultados referentes a los estadísticos y sus propiedades,
los test de hipótesis nula, tanto paramétricos como no 
paramétricos y la estadística bayesiana.\\
	Los conceptos informáticos se han extraído del área
del aprendizaje automático, y se corresponden con el 
problema de clasificación, las medidas para evaluar el 
rendimiento o los algoritmos de clasificación.

\section*{Contenido de la memoria}\addcontentsline{toc}{section}{Contenido de la memoria}
	En la primera parte de la memoria
(\ref{part:matematicas}) se presenta en el 
capítulo~\ref{chapter:Intro} una introducción a la evaluación 
de algoritmos de aprendizaje automático, a las medidas 
utilizadas para ello y la notación usada. En los 
sucesivos capítulos se presenta el contenido matemático
relativo a los test paramétricos 
(Cap.~\ref{chapter:test_parametricos}), los test no
paramétricos y basados en permutaciones 
(Cap.~\ref{chapter:test_no_parametricos}) y
los test bayesianos (Cap.~\ref{chapter:test_bayesianos}).\\
	En la segunda parte de la memoria 
(\ref{part:informatica}) se hace en el capítulo~
\ref{chapter:revision} una revisión del estado del arte
en la aplicación de test estadísticos al problema que 
nos ocupa. En el Capítulo~\ref{chapter:software} se
realiza una descripción de varios paquetes software 
existentes para la realización de test estadísticos, los
cuales se integran en un paquete de \texttt{R} que funciona
como \textit{wrapper}. Finalmente,
en el capítulo~\ref{chapter:experimentos} se realiza un
caso práctico donde se aplican los test estadísticos
disponibles y se compara su potencia.

\section*{Principales fuentes utilizadas}\addcontentsline{toc}{section}{Fuentes utilizadas}

	Las principales fuentes utilizadas han sido los libros de
Sheskin
(\cite{sheskin2003handbook}), Japkowicz
(\cite{DBLP:books/cu/Japkowicz2011})
y Gibbons y Chakraborti (\cite{DBLP:reference/stat/GibbonsC11}) 
para el desarrollo matemático
y la descripción de los test estadísticos paramétricos y
no paramétricos. Los métodos paramétricos están basados 
también en los trabajos de Alpaydin (\cite{Alpaydin98combined5},
\cite{Alpaydin:2010:IML:1734076}) y Ditterich
 (\cite{dietterich1998approximate}).
Las metodologías no paramétricas se
han complementado principalmente con los artículos de 
Dem\v{s}ar (\cite{DBLP:journals/jmlr/Demsar06}) y 
García (\cite{garcia2008extension}, 
\cite{DBLP:journals/isci/GarciaFLH10}). Para los test basados
en permutaciones se ha utilizado fundamentalmente el 
libro de Pesarin y Salmaso (\cite{pesarin2010permutation}).
Los test bayesianos han sido desarrollados usando 
principalmente el tutorial de Benavoli \textit{et. al} 
(\cite{DBLP:journals/corr/BenavoliCDZ16}).\\
	Para la parte de informática, se han utilizado la
documentación disponible de los paquetes integrados
(\cite{2015arXiv150104222D}, \cite{scmamp}, 
\cite{DBLP:conf/icml/BenavoliCMZ15}).

\chapter*{Objetivos}\addcontentsline{toc}{chapter}{Objetivos}

	El primer objetivo planteado fue realizar la 
recopilación, el estudio bibliográfico y recoger la 
taxonomía de las técnicas estadísticas que se aplican
en estudios experimentales en el área del aprendizaje
automático. La realización de este objetivo se realiza
en la primera parte de la memoria (\ref{part:informatica}) y
en la sección dedicada al estudio del estado del arte 
(\ref{chapter:revision}), incluyendo también entre los
test descritos algunos relacionados con comparación a
pesar de no ser estrictamente ésta su finalidad
como los test de aleatoriedad o los de bondad del 
ajuste.\\
	El segundo fue el diseño, desarrollo
e implementación de una biblioteca que unificase los test
estadísticos más potentes y utilizados, objetivo 
realizado basándonos en bibliotecas ya creadas para 
test específicos y que han sido unificadas en una biblioteca
de \texttt{R} (\ref{chapter:software}).\\
	El último objetivo planteado fue la evaluación y el 
análisis experimental de las diferentes técnicas estadísticas 
desarrolladas, para lo que se han obtenido datos 
experimentales del rendimiento de varios algoritmos en 
múltiples conjuntos de  datos sobre los que se aplican los 
test descritos e implementados para realizar una comparación 
acerca de su potencia y la información extraída de ellos
(\ref{chapter:experimentos}).



	


%
\part{Desarrollo matemático}
\label{part:matematicas}

\chapter{Notación e introducción a la evaluación de algoritmos de aprendizaje}
\label{chapter:Intro}

	En este capítulo se hace un breve repaso de los 
conceptos estadísticos necesarios para comprender el 
contenido de la memoria así como presentar la notación,
una introducción con las medidas de rendimiento y 
diferentes formas de utilizar los datos para 
obtener un valor más preciso del rendimiento.

\section*{Notación}
	
\begin{definicion}[Variable aleatoria]
	Función de conjuntos cuyo dominio es los elementos de un 
espacio muestral sobre el cual se ha definido una función de
probabilidad y cuyo rango es $\mathbb{R}$.\\
	$X$ es una variable aleatoria (v.a.) si para $x \in 
\mathbb{R}$ existe una probabilidad de que el valor tomado 
por la variable aleatoria sea menor o igual que $x$, es
decir, $P(X \leq x) = F_X (x)$, llamada función de
distribución de probabilidad (\textit{cumulative distribution
function}, cdf) de $X$.	
\end{definicion}

	Cualquier función de distribución, $F_X(x)$, de una v.a.
$X$ cumple las siguientes propiedades:
 
\begin{enumerate}
	\item $F_X$ es no decreciente: 
			$F_X(x_1) \leq F_X(x_2) \ \forall x_1 \leq x_2$.
	\item $\underset{x \rightarrow -\infty}{\lim} F_X(x) =
			 0$,
			$\underset{x \rightarrow \infty}{\lim} F_X(x) =
			 1$
	\item $F_X(x)$ es continua por la derecha: 
		$\underset{\varepsilon \rightarrow 0^+}{\lim} 
		F_X(x+\varepsilon) = F_X(x)$
\end{enumerate}
	
\begin{definicion}[Función de densidad]
	Se define la función de densidad como la derivada de 
$F_X(x)$, $f_X(x)$. Para $X$ continua:
	\[ 
		F_X(x) = \int_{-\infty}^x f_X(t) dt; \;
		f_X(x) = \frac{d}{dx}F_X(x) = F_X'(x) \geq 0; \;
		\int_{-\infty}^{\infty} f_X(x) dx = 1 
	\]
\end{definicion}
	
\begin{definicion}[Función de masa]
	Se define la función de masa de probabilidad  
(\textit{probability mass function}, pmf) de una v.a. 
\textbf{discreta}, es decir, que sólo toma un número contable 
de valores como
	\[ 
		f_X(x) = P(X=x) = 
		F_X(x) - 
			\lim_{\varepsilon \rightarrow 0^+} 
				F_X(X-\varepsilon).
	\]
\end{definicion}

\begin{definicion}[Muestra aleatoria]
	Llamamos muestra aleatoria de una v.a. $X$ a un conjunto 
de $n$ v.a., $\mathbf{X} = (X_1, \dots, X_n)$, si son independientes e 
idénticamente distribuidas (i.i.d.), con lo que su 
distribución de probabilidad conjunta es
	\[ 
		f_{\mathbf{X}}(x_1, \dots, x_n) =
		f_{X_1, \dots, X_n}(x_1, \dots, x_n) =
		\prod\limits_{i=1}^n f_X(x_i).
	\]
\end{definicion}
	
\begin{definicion}[Momento]
	Es un parámetro de la población. El momento $k$-ésimo de 
$X$ es $\mu_k' = \mathbb{E}[X^k]$. La media es el momento de 
primer orden, $\mu_1' = \mathbb{E}[X] = \mu$. El momento 
$k$-ésimo centrado es  $\mu_k = \mathbb{E}[(X - \mu)^k]$.
Notamos por $\var$ a la varianza, es decir, el momento 
centrado de segundo orden $\var(X) = \mathbb{E}[(X - \mu)^2]$. 
\end{definicion}

\begin{definicion}[Función generadora de momentos]
	Notamos por $M_{X}(t)$ a la función $\mathbb{E}[
\exp(tX)]$. 
\end{definicion}

	Sea $\hat{\theta}_n = u(X_1, \dots, X_n)$ un estimador de 
un parámetro $\theta$. Incluimos unas propiedades deseables 
de $\hat{\theta}_n$:
	
	\begin{enumerate}
	\item \textit{Insesgadez}: 
			$\mathbb{E}[\hat{\theta}_n] = \theta$ para todo $\theta$.
	\item \textit{Suficiencia}: Podemos escribir
			$f_{X_1, \dots, X_n}(x_1, \dots, x_n; \theta)$ 
			como producto de dos funciones $f_{X_1, \dots, 
			X_n}(x_1, \dots, x_n; \theta) = g(\hat{\theta}_n; 
			\theta) H((x_1, \dots, x_n)$ tal que $H(x_1, 
			\dots, x_n)$ no depende de $\theta$.
	\item \textit{Consistencia}
		\[ 
			\lim_{n \rightarrow \infty} 
				P(|\hat{\theta}_n - \theta| < \varepsilon) =
				 0 \quad 
				 \forall \varepsilon > 0 
		\]
	\begin{enumerate}
		\item Si $\hat{\theta}_n$ es un estimador insesgado 
			de $\theta$ y $\underset{n \rightarrow \infty}
			{\lim} \var(\hat{\theta}_n) = 0$, entonces $
			\hat{\theta}_n$ es un estimador consistente por la 
			desigualdad de Chebyshev.
		\item $\hat{\theta}_n$ es un estimador consistente de 
			$\theta$ si la distribución límite es la 
			distribución degenerada con probabilidad $1$ en 
			$\theta$.
	\end{enumerate}
			
	\item \textit{Mínimo error cuadrático:} 	
			$\mathbb{E}[(\hat{\theta}_n - \theta)^2] \leq
			 \mathbb{E}[(\hat{\theta}_n^* - \theta)^2]$ para cualquier 
			 estimador $\hat{\theta}_n^*$.
	\item \textit{Mínima varianza:} $\var(\hat{\theta}_n)
			\leq \var(\hat{\theta}^*_n)$ para cualquier 
			estimador $\hat{\theta}_n^*$, siendo ambos
			insesgados.
	\end{enumerate}
	
\begin{definicion}[Función de verosimilitud]
	La función de verosimilitud (\textit{likelihood 
function}) de una muestra aleatoria de tamaño $n$ de la 
población $f_X(x;\theta)$ es la probabilidad conjunta de las 
muestras tomadas como función de $\theta$. Esto es:
	\[ L(x_1, \dots, x_n; \theta) = 
		\prod\limits_{i=1}^n f_X(x_i;\theta)	\]
\end{definicion}

	Un \textbf{estimador máximo verosímil} (MLE) de $\theta$ 
es un valor $\bar{\theta}$ tal que 
	\[ 
		L(x_1, \dots, x_n; \hat{\theta}) \geq 
			L(x_1, \dots, x_n; \theta) \; \forall \theta 	
	\]
	
	La relevancia de este estimador consiste en que, para 
unas ciertas condiciones de regularidad, un estimador máximo 
verosímil es suficiente, consistente y asintóticamente 
insesgado, con varianza mínima y con distribución normal.
	
\begin{definicion}[Hipótesis estadística]
	Es una afirmación sobre la la función de probabilidad de 
una o más v.a. o una afirmación sobre las poblaciones de las 
cuales se han obtenido una o más muestras aleatorias. La 
\textbf{hipótesis nula}, $H_0$ es la hipótesis sobre la que 
se realizará un test. La \textbf{hipótesis alternativa}, 
$H_1$ es la que conclusión alcanzada si se rechaza la 
hipótesis nula.
\end{definicion}

\begin{definicion}[Región crítica]
	Llamamos región crítica o región de rechazo $R$ para un 
test al conjunto de valores tomados por el test que conducen 
a rechazar la hipótesis nula. Llamamos \textbf{valores 
críticos} a los extremos de $R$ y los notaremos por $\tau_
\alpha$
\end{definicion}

\begin{definicion}[Tipos de error]\textit{}
	\begin{description}
	\item[Error de tipo I] La hipótesis nula es rechazada 
		siendo cierta.
	\item[Error de tipo II] La hipótesis nula no es rechazada 
		siendo falsa.
	\end{description}
\end{definicion}

	Siendo $T$ un test estadístico con hipótesis $H_0: \theta 
\in \omega, \ H_1: \theta \in \Omega \setminus \omega$, los 
errores de tipo I y II tienen probabilidad
	\[ 
	\alpha(\theta) = P(T \in R | \theta \in \omega); \quad
	\beta(\theta) = 
		P(T \not\in R | 
				\theta \in \Omega \setminus \omega)
	\]
	respectivamente.

\begin{definicion}[Tamaño del test]
	Se define el tamaño del test como $\underset{\theta \in 
	\omega}{\sup}\ \alpha(\theta)$.
\end{definicion}

\begin{definicion}[Potencia del test]
	Se define la potencia del test como la probabilidad de 
que el test conduzca a un rechazo de $H_0$: $Pw(\theta) = P(T 
\in R)$. Esta medida nos interesa cuando debemos rechazar la 
hipótesis nula, con lo que calculamos $Pw(\theta) = P(T \in R 
| \theta \in \Omega \setminus \omega) = 1 - \beta(\theta)$. 
\end{definicion}
	
	Diremos que un test es \textbf{más potente} para una 
hipótesis alternativa concreta si ningún test del mismo 
tamaño tiene mayor potencia contra la misma hipótesis 
alternativa.\\
	A continuación definimos un enfoque distinto,
especialmente relevante en los test no paramétricos. 
	
\begin{definicion}[$p$-valor]
	Probabilidad de obtener, siendo cierta la hipótesis nula 
$H_0$, una muestra aleatoria tan extrema como la muestra
aleatoria observada.
\end{definicion}

\begin{definicion}[Consistencia]
	Diremos que un test es consistente para una hipótesis 
alternativa $H_1$ si la potencia del test se aproxima a 1 
conforme $n \rightarrow \infty$, siendo $n$ el tamaño de la 
muestra.
\end{definicion}

	Notaremos como $\mathbb{I}[\textit{condición}]$ a la 
función indicadora, que toma valor $1$ si \textit{condición} es
cierta y 0 si no.


\section*{Medidas del rendimiento}
	
	El rendimiento de un algoritmo sobre un problema puede 
depender del criterio que los investigadores deseen optimizar 
con lo que el concepto de la \textit{calidad} de un algoritmo 
depende según el problema al que se aplique. Por tanto, 
consideraremos el problema como un problema multiobjetivo. 
Una posible forma de abordarlo es condensar varias medidas en 
una única con la que posteriormente realizar las 
comparaciones necesarias a través de test estadísticos. Se 
presenta a continuación un resumen de las medidas comúnmente 
utilizadas en el campo del aprendizaje automático tanto para 
problemas de clasificación como de regresión, indicando 
ventajas e inconvenientes que nos ayuden a decidir aquella 
medida que mejor se ajuste a nuestro problema.
	
\subsection*{Medidas para problemas de clasificación}

	Las medidas utilizadas habitualmente para problemas de 
clasificación están basadas en la matriz de confusión 
(~\ref{tab:matrizconfusion}). Haremos un listado de aquellas 
más utilizadas sin entrar en las medidas diseñadas \textit{ad 
hoc} para problemas concretos. Suponemos que nos encontramos 
frente a un problema de clasificación binaria, aunque es 
posible extender la matriz de confusión y las medidas para 
problemas de clasificación con un mayor número de clases.
Notaremos por $x, \bar{x}$ la clase positiva y la negativa
respectivamente de nuestro problema, mientras que $y,\bar{y}$
se corresponden con las asignaciones a la clase positiva
y negativa respectivamente. Por tanto, notaremos por $TP$ a
los verdaderos positivos, por $TN$ a los verdaderos negativos,
$FP$ a los falsos positivos y $FN$ a los falsos negativos.

\begin{table}[H]
\centering
\begin{tabular}{cccc}
\cline{1-3}
\multicolumn{1}{|c|}{}           & \multicolumn{2}{c|}{Predicción del clasificador}      &                 \\ \cline{1-3}
\multicolumn{1}{|c|}{Clase real} & \multicolumn{1}{c|}{$y$}    & \multicolumn{1}{c|}{$\bar{y}$}    &                 \\ \cline{1-3}
\multicolumn{1}{|c|}{$x$}          & \multicolumn{1}{c|}{$TP$} & \multicolumn{1}{c|}{$FN$} & $N^+ = TP + FN$ \\ \cline{1-3}
\multicolumn{1}{|c|}{$\bar{x}$}          & \multicolumn{1}{c|}{$FP$} & \multicolumn{1}{c|}{$TN$} & $N^- = FP + TN$ \\ \cline{1-3}
                                & $\hat{N}^+ = TP + FP$     & $\hat{N}^+ = FN + TN$     &   $N = N^- + N^+$             
\end{tabular}
\caption{Matriz de confusión}
\label{tab:matrizconfusion}
\end{table} 

En función de esta matriz se definen las siguientes
medidas.
	
\begin{description}
	\item[Error de clasificación ($\varepsilon$),] 
		$\varepsilon = \frac{FP+FN}{N}$
	\item[Exactitud o \textit{Accuracy}(Acc),] 
		$Acc=\frac{TP+TN}{N}$
	\item[Exhaustividad o \textit{Recall}(r),] $r=\frac{TP}{N^+}$
	\item[Especificidad o \textit{Specificity}(sp),] 
		$sp=\frac{TN}{N^-}$
	\item[Precisión (pr),] 
		$pr=\frac{TP}{\hat{N}^+}$
	\item[Tasa falsos positivos(fpr)] $fpr=1-sp$
	\item[Tasa falsos negativos(fnr)] $fnr=1-r$
	\item[Valor negativo predicho(npv)] 
		$npv=\frac{TN}{\hat{N}^-}$
	\item[Estadístico kappa ($\kappa$)] El propósito de esta 
		medida es medir el grado de acuerdo entre dos 
		observadores, o, en nuestro contexto, compensar los 
		valores aleatorios además de minimizar el efecto de 
		un dispar número de ejemplos en cada clase. 
	\[ 
	\kappa = \frac{N(TP + TN) - N^+\hat{N}^+ - N^-\hat{N}^- }
				  {N^2 - N^+\hat{N}^+ - N^-\hat{N}^- }
	\]
\end{description}
	
	Las medidas de exactitud y el error de 
clasificación son comúnmente usadas debido a que representa 
un concepto claro y son fáciles de calcular.  Es necesario 
indicar que estas medidas son aconsejables cuando las dos 
clases se dan con una frecuencia similar y son igualmente 
relevantes para nuestro problema. Otras medidas como 
\textit{recall}, \textit{specificity} o \textit{precision} 
son más frecuentes en dominios como el de recuperación de 
información y aquellos donde o bien una clase es más 
relevante, o sólo existen unos pocos individuos de 
una clase.\\
	Para encontrar un equilibrio entre la detección de 
elementos de la muestra de ambas clases se han propuesto en 
la literatura (por ejemplo por Santafé, 
\cite{DBLP:journals/air/SantafeIL15}) algunas aproximaciones 
basadas en las anteriores medidas:
	
	\begin{description}
		\item[Compensación entre exhaustividad y 
especificidad.] El método más popular para encontrar 
este equilibrio es el análisis ROC, que consta de una 
representación de $r$ con respecto a $fpr$. Una curva ROC 
puede resumirse en el valor AUC (\textit{area under the ROC 
curve})
		\item[Compensación entre exhaustividad y 
precisión.] Depende del parámetro $\beta$ dado por el 
usuario :
		\[ F_\beta = \frac{(\beta^2+1) pr \cdot r}
						  {\beta^2 pr + r}	\]
	\end{description}

	Cuando en un problema de clasificación supervisada usamos 
un modelo probabilístico obtenemos la probabilidad de 
pertenencia a una de las clases, con lo que podemos realizar 
la asignación variando el umbral. Existen medidas que tratan 
de recoger el grado de proximidad entre esta asignación de 
probabilidad y la distribución real. La más conocida es la 
medida Brier ($BS$),
	\[ BS = \frac{1}{N} \sum\limits_{i=1}^N 
		\sum\limits_{j \in \{c^+,c^-\}}
			\left( P \left( C=j|\mathbf{x}^{(i)} \right) -
			\delta_{j,c^{(i)}} \right) ^2,\]
	donde $P\left(C=j|\mathbf{x}^{(i)}\right)$ es la 
probabilidad asignada por el clasificador a la instancia 
$i$-ésima para la clase $j$, $\delta$ es la delta de 
Kronecker y $c^{(i)}$ es la clase de la instancia 
$i$-ésima.\\
	Hasta ahora las medidas le han dado el mismo peso a los 
errores cometidos en ambos sentidos (asignar a la clase 
positiva una instancia negativa o viceversa), sin embargo se 
pueden incluir \textit{funciones de pérdida} ($\mathcal{L}$) 
según el coste relativo de cada tipo de fallo. Se 
puede considerar que en las medidas anteriores se ha usado 
la función de pérdida 
$\mathcal{L}(x,y) = 1 - \delta_{x,y} $, donde $x$ es la clase 
real e $y$ la clase asignada.
	
\subsubsection*{Medidas de interpretabilidad} 
	En algunos modelos y clasificadores se obtienen reglas 
para la clasificación de nuevos ejemplos. Por ello se valora 
positivamente que estas reglas sean entendibles para el 
usuario. 
	\begin{description}
	\item[Tamaño] Número de reglas que componen el modelo. 
		A menor tamaño, mayor interpretabilidad. 
		$\textit{Tamaño} = n_R$
	\item[Número de antecedentes (ANT)] Para $R_i$ una regla 
		de la forma $\textit{Condi-}$ ${ciones } \rightarrow 
		\textit{ Clase}$ donde las condiciones tienen la 
		forma $(Antecedente_1 \wedge \dots \wedge 
		Antecedente_k)$, definimos $Ant(R_i) = k$ y entonces 
		\[ANT = \frac{1}{n_R} 
				\sum\limits_{i=1}^{n_R} Ant(R_i).\]
	\end{description}


\subsubsection*{Medidas gráficas}

	Prati \textit{et al.} \cite{DBLP:journals/tkde/PratiBM11} afirman que en ocasiones es 
recomendable utilizar una representación gráfica del 
rendimiento ya que evitamos así pérdida de información al 
condensar toda la información en un escalar. En este artículo 
se presenta un conjunto de representaciones gráficas según el 
tipo de variable predicha: clases, ránquin o probabilidad de 
pertenencia a una clase.
	
\paragraph{Predicción de clases}
	\begin{description}
		\item[Gráficas ROC] Como se ha mencionado 
		anteriormente, es una representación del la medida 
		\textit{tpr} frente a \textit{fpr}. Un clasificador 
		será mejor considerado cuanto más arriba y a la 
		izquierda se sitúe en el gráfico. 
		\item[Líneas de coste] Se trata de una modificación 
		de las gráficas ROC para incluir el coste asociado a 
		cada tipo de fallo. Se define el coste esperado para 
		un clasificador como 
		\[ EC = \sum\limits_{X \in \{x, \bar{x}\}} 
		 		\sum\limits_{X \in \{x, \bar{x}\}} 
		 			P(X,Y) \mathcal{L}(X,Y),
		 \]
con $P(X,Y)$ la probabilidad de clasificar como $Y$ un
elemento de la clase $X$ y $\mathcal{L}(X,Y)$ la
función de coste asociado a considerar $Y$ un elemento
de la clase $X$.
Entonces, se definen $NormEC = \frac{EC}{maxEC}$, $PC = 
\frac{p(x) \mathcal{L}(\bar{x},y)}{maxEC}$, donde $maxEC 
= p(x)\mathcal{L}(\bar{x},y)+p(\bar{x})\mathcal{L}(x,
\bar{y})$.  La gráfica representa la $PC$ frente a 
$NormEC$. Ahora cada clasificador, representado por un 
punto en la gráfica ROC, es una línea. El clasificador 
representado por el punto $(0,0)$ en la gráfica ROC (asigna a 
todos los puntos la clase negativa) ahora se representa 
mediante la recta $NormEC=PC$ y el punto $(1,1)$ mediante la 
recta $NormEC= 1-PC$. Dependiendo del coste relativo de la 
clase positiva ($PC$), será el mejor clasificador aquel con 
menor $NormEC$.
	\end{description}

\paragraph{Predicción de \textit{ranking}} 	
	\begin{description}
		\item[Curvas ROC] Se traslada el problema a uno de 
		clasificación mediante la asignación a la clase 
		positiva del \textit{top} $n\%$. Se varía el 
		porcentaje del $0\%$ al $100\%$ para obtener la curva 
		ROC.
		\item[Curvas de coste] Se trata de la transformación 
		análoga a las curvas ROC para las líneas de coste.
		\item[Curvas \textit{precision-recall}] Usadas 
		habitualmente en el ámbito de la recuperación de 
		información. Representan $precision$ frente a 
		$recall$.
		\item[Gráfica \textit{lift}] Similar a la gráfica 
		ROC, representa $tpr$ frente a la probabilidad de 
		asignar a un individuo la clase positiva. 
		\item[Gráfica ROI] Similar a las gráficas 
		\textit{lift}, pero representando en el eje $Y$ el 
		beneficio esperado 
		\[	TEP = N \sum\limits_{X \in \{x,\bar{x}\}}
					\sum\limits_{Y \in \{y,\bar{y}\}}
						p(X,Y) \mathcal{L}(X,Y) p(X) \]
	\end{description}

\paragraph{Predicción de probabilidad de pertenencia a una clase} 	
	\begin{description}
		\item[\textit{Reliability diagram}] Contiene dos
		gráficas. Una es $p(x|j)$ según $j$, que mide el 
		grado en que concuerdan las predicciones con la
		frecuencia de los casos positivos. La segunda gráfica 
		se corresponde con $p(y_j)$ según $j$, identificada 
		con la confianza en la predicción.
		\item[Diagrama de atributos] Para evaluar las mejoras 
		relativas de un sistema de predicción con respecto a 
		uno de referencia se define la $BSS$ (\textit{Brier 
		skill score})
			\[BSS = 1 - \frac{BS}{BS_{ref}} \]
		El diagrama de atributos es un refinamiento del
		\textit{reliability diagram}, al que se le añaden 
		tres líneas: 
			\begin{description}
			\item[Línea de no resolución:] línea horizontal 
			donde $p(x|j)=p(x)$.
			\item[Línea de máxima resolución:] línea vertical 
			que intersecta a la diagonal en el mismo punto 
			que la línea de no resolución.
			\item[Clasificador de referencia:] Bisectriz 
			entre la diagonal principal y la línea de no 
			resolución. Representa los puntos cuya precisión 
			(\textit{accuracy}) es menor que en el modelo de 
			referencia.
			\end{description}
		\item[\textit{Discrimination diagram}] Gráfica de 
		$P(Y|X)$ según $y_j$. Contiene las curvas de 
		$p(y_j|x)$ y de $p(y_j|\bar{x})$. En el mejor de los 
		casos, estas dos curvas no se solaparían y podríamos 
		saber la clase según $y_j$.
	\end{description}
	
	
\subsection*{Medidas para problemas de regresión}

	Las medidas para un problema de regresión están basadas 
en la distancia entre el valor real de la instancia y el 
valor predicho. Consideraremos que calculamos el error de la 
predicción del algoritmo $f$ aplicado a los datos $x_1, 
\dots, x_N$ con valor $y_1, \dots, y_N$ respectivamente.
	
\begin{description}
	\item[\textit{Mean squared error} (MSE)]: Error 
	cuadrático medio
		\[ MSE(f) = \frac{1}{N} \sum\limits_{i=1}^N
								(y_i - f(x_i))^2	\]
	\item[\textit{Root mean squared error} (RMSE)]: Raíz del 
	error cuadrático medio
		\[ RMSE(f) = \sqrt{\frac{1}{N} 
		\sum\limits_{i=1}^N (y_i - f(x_i))^2}	\]
		
	\item[\textit{Mean absolute error} (MAE)]: Error 
	absoluto medio
		\[ MAE(f) = \frac{1}{N} \sum\limits_{i=1}^N
								|y_i - f(x_i)|	\]
\end{description}

	Las medidas $MSE$ y $RMSE$ son más utilizadas debido a la 
penalización de los errores mayores. 

\section{Validación cruzada y remuestreo}
\label{sec:CV}
	
	Para la repetición de los experimentos, necesitamos un 
número de conjuntos de entrenamiento y validación del 
conjunto de datos disponible. Si el conjunto es 
suficientemente grande, podemos dividir aleatoriamente el 
conjunto en $K$ partes, y de cada parte seleccionar la mitad 
para entrenamiento y la otra mitad para contrastar. $K$ suele ser 
10 ó 30. Sin embargo, los conjuntos de datos no suelen ser 
tan grandes. Por tanto, lo que habitualmente se realiza es 
utilizar varias veces los mismos datos pero de formas 
distintas, lo que se conoce como \textbf{validación cruzada} 
(\textit{cross-validation}, CV). Otro concepto a tener en 
cuenta es el de \textbf{estratificación}. Consiste en 
mantener la proporción de las diferentes clases en cada una 
de las particiones realizadas en la base de datos. 
	
\paragraph{Validación cruzada con $K$-fold} Se divide el 
conjunto de datos $\mathcal{X}$ en $K$ partes aleatoriamente, 
$\mathcal{X} = \mathcal{X}_1 \cup \dots \cup \mathcal{X}_K$. 
Para formar cada par de datos de entrenamiento y validación, 
$(\mathcal{T}, \mathcal{V})$, se mantiene una de las $K$ 
partes y se unen las demás partes como conjunto de 
entrenamiento
\begin{align*}
	\mathcal{V}_i &= \mathcal{X}_i \\
	\mathcal{T}_i &= \underset{j=1,\dots,K; j\neq i}
							\bigcup \mathcal{X}_j
\end{align*} 
	El problema de este método es que el conjunto de 
validación es pequeño (lo que lleva a una mayor variabilidad 
en el error). Además, los conjuntos de entrenamiento se 
solapan considerablemente. $K$ es habitualmente 10 ó 30. 
Conforme $K$ crece, se incrementa la robustez al afectar 
menos cada dato concreto. El caso $K = n-1$, se conoce como 
\textit{leave-one-out}. 

\paragraph{$5 \times 2$ CV} Otra propuesta consiste en 
dividir $\mathcal{X}$ en dos partes, utilizar una de estas 
partes como datos de entrenamiento y la otra como test e 
invertir los roles. Para obtener el segundo \textit{fold}, se 
realiza otra partición aleatoria de $\mathcal{X}$ y se repite 
el proceso. Lo habitual es que se realice cinco veces, 
obteniendo así diez parejas $(\mathcal{T}, \mathcal{V})$. 
Aunque podría realizarse más veces, al realizarse más veces 
los conjuntos comparten muchas instancias y por lo tanto los 
errores son muy dependientes entre sí, con lo que no se 
aporta nueva información. Si se realiza menos de cinco veces, 
se obtienen pocas muestras y es más complejo comprobar las 
hipótesis.

\paragraph{\textit{Bootstrapping}} Para generar varias 
muestras a partir de una única se seleccionan instancias de 
la original con reemplazamiento. Puede solaparse más que con 
CV, con lo que las estimaciones resultan más dependientes. 
Por ello, se recomienda especialmente para conjuntos de datos 
muy pequeños en los que se necesiten mayores muestras. 

	Al realizar \textit{bootstrap}, se seleccionan $N$ 
instancias de un conjunto de tamaño $N$ con reemplazamiento. 
El conjunto original se usa como conjunto de validación. La 
probabilidad de no seleccionar una instancia es $1-1/N$, con 
lo que la probabilidad de no seleccionarla tras $N$ 
elecciones es $\left(1- \frac{1}{N}\right)^N \approx e^{-1} 
\approx 0.368$. Esto significa que el conjunto de 
entrenamiento contiene aproximadamente un $63.2\%$ de los 
datos. Para obtener una mejor estimación del error, se 
propone repetir el proceso y observar el comportamiento 
medio.
	
	
\chapter{Test paramétricos}
\label{chapter:test_parametricos}

	Los test paramétricos se basan en la suposición de que la
muestra pertenece a una distribución que sigue un modelo 
conocido (frecuentemente, un modelo gaussiano). La ventaja de 
este enfoque es que la definición del modelo atiende a un 
pequeño número de parámetros (son los estadísticos 
suficientes). Estimando estos parámetros a partir de la 
muestra, se conoce la distribución. 
	
\section{Test binomial para una muestra}

	Con este test se pretende comprobar si en una población
compuesta por dos categorías de la que se extrae la muestra, 
la proporción de observaciones de una de las categorías es 
igual a un valor concreto. En nuestro contexto, consideramos 
$\theta$ la probabilidad del clasificador de cometer un error 
y queremos realizar una hipótesis sobre el valor de $\theta$. 
Supongamos que tenemos una muestra de tamaño $N$, 
$\{(x_i, y_i): i = 1, \dots,N \}$ con $y_i$ la categoría y
consideraremos la variable aleatoria $X$ que denota el número
de errores cometidos por el clasificador $f$,
	\[ 
		X = \sum\limits_{i=1}^N 
				\mathbb{I}[f(x_i) \neq y_i].
	\]
	
	El test a realizar consiste en comprobar si la probabilidad de 
cometer un error $\theta$ es igual a un valor $\theta_0$
dado: $H_0: \theta = \theta_0;\ H_1: \theta \neq \theta_0$.
La probabilidad de cometer $j$ errores de $N$ supuesta $H_0$
será $P[X = j] = {N \choose j} \theta^j (1-\theta)^{N-j}$. Se 
rechazará la hipótesis nula si el $p$-valor observado 
es menor que el valor $\alpha/2$ (al considerar la hipótesis 
alternativa no direccional).\\
	Cuando el tamaño de la muestra es grande el test 
estadístico para la binomial puede ser aproximado con la 
distribución $\chi^2$ con un grado de libertad. Hay varias 
opiniones con respecto al tamaño mínimo para realizar
esta aproximación. Se pueden tomar referencia que 
$N \theta_0$ y $N (1 - \theta_0)$ sean mayores que 5. El 
estadístico para realizar el test usando esta 
aproximación será:
	\[
		z = \frac{X/N - \theta_0}
				{\sqrt{\frac{\theta_0(1 - \theta_0)}
							{N}}}
	\]
	
\section{$t$-Test para muestras apareadas}

	Nos ocupamos aquí de esta versión del $t$-test para 
muestras independientes apareadas con varianzas 
desconocidas por tener una mayor relevancia en la evaluación 
de algoritmos de aprendizaje automático. Suponiendo que 
disponemos de los datos del rendimiento de dos algoritmos, 
realizamos el test para comprobar si éstos se corresponden 
con los propios de dos poblaciones con diferentes medias. 
Para estimar el valor de las medias de las poblaciones, 
notaremos por $\bar{X}_1, \bar{X}_2$ a las medias de las dos 
muestras. Las condiciones para realizar este test son:
	
	\begin{itemize}
	\item Cada muestra ha sido seleccionada aleatoriamente 
		de la población a la que representa.
	\item La distribución de los datos de las poblaciones de 
		las que se extraen las muestras son normales.
	\item Homogeneidad de las varianzas: La varianza de las 
		dos poblaciones son iguales ($\sigma_1 = \sigma_2$).
	\end{itemize}
	
	Nótese que si la varianza es conocida, es más apropiado 
usar el $z$-test.\\
	La hipótesis nula será por tanto $H_0: \mu_1 = \mu_2$.
El test se realiza utilizando el estadístico $t$:
\begin{align*}
	t 	= \frac{\bar{d}}
			{\frac{\bar{\sigma}_d}{\sqrt{N}}}
		= \frac{\bar{X}_1 - \bar{X}_2}
			{\frac{\bar{\sigma}_d}{\sqrt{N}}},
\end{align*}
	
	donde $N$ es el número de datos de la muestra, 
$\bar{X}_i,\ i=1,2$  es la media $\frac{1}{N}
\sum\limits_{j=1}^{N} X_{ij}$ de los  valores obtenidos en 
cada muestra, $\bar{d}$ es la diferencia de las medias y 
$\bar{\sigma}_d$ es el estimador para la desviación típica
de las diferencias:
	\[
		\bar{\sigma}_d =
			\sqrt{
			\frac{\sum\limits_{i=1}^{N} (d_i - \bar{d})^2}
				{N - 1}
			}.
	\]
	
	El valor $t$ obtenido se comparará con el valor crítico
para la distribución $t$ de Student con $N - 1$ 
grados de libertad.\\
	Una versión más general de este test es el $t$-test de
Welch para muestras no apareadas de tamaños $N_1, N_2$
cuyo estadístico es
	\[
		t = \frac{\bar{X}_1 - \bar{X}_2}
				{\sqrt{ \frac{\sigma_1^2}{N_1} +
						\frac{\sigma_2^2}{N_2}}},
	\]
	con $\sigma_i^2,\ i=1,2$ la varianza de la muestra.
	
\paragraph{Tamaño del efecto} El $t$-test determina si las 
diferencias observadas son significativas, pero no si 
estas diferencias tienen importancia práctica, es decir, 
se observa un efecto, una diferencia entre los 
clasificadores, pero no se cuantifica esta diferencia
(crítica que se le hace desde la estadística bayesiana,
como veremos). El tamaño del efecto para este test se
suele realizar con el estadístico $d$ de Cohen:
	\[
		d_{Cohen} = \frac{\bar{X_1}-\bar{X_2}}
					{\sigma_p},
	\]
	con $\sigma_p = \sqrt{\frac{\sigma_1^2 + \sigma_2^2}{2}}$
la raíz de la media de las varianzas de las respectivas muestras.
La interpretación del estadístico propuesta por Cohen es:
	\begin{itemize}
	\item $d_{Cohen}$ en torno a $0.2, 0.3$ denota un 
		efecto pequeño, posiblemente significativo.
	\item $d_{Cohen}$ sobre $0.5$ significa un efecto medio,
		observable.
	\item $d_{Cohen}$ de $0.8$ significa un gran tamaño del 
		efecto.
	\end{itemize}
	
		
\section{Análisis de la varianza}

	Estos test resultan insuficientes a la hora de comparar 
múltiples clasificadores. Aunque puede usarse el $t$-test para 
realizar repetidas comparaciones entre distintos 
clasificadores, esto plantea el problema de que son 
necesarios numerosos test para realizar todas las 
combinaciones y por tanto la información que se obtiene es
más difícil de analizar. Otro problema asociado es el 
incremento de la probabilidad de cometer un error de tipo I. 
Presentamos en esta sección la versión paramétrica, ANOVA.\\
	En este test, al igual que en el $t$-test, la hipótesis 
hace referencia a las medias de las poblaciones de las que 
se extraen las muestras, sin embargo, la hipótesis nula es
	\[
		H_0 : \mu_1 = \mu_2 = \dots = \mu_k,
	\]
	frente a la hipótesis alternativa
	\[
		H_1 : \mu_r \neq \mu_s \text{ para algún par }
			(r, s), r \neq s 
	\]
	
	El test ANOVA puede tratar con tres tipos de varianzas
si los datos están clasificados en grupos (como es nuestro
caso, donde los datos se podrán agrupar bien por los errores
de todos los clasificadores en una base de datos, o bien por
los errores de un clasificador en todas las conjuntos de datos):
	\begin{enumerate}
	\item variación en un grupo,
	\item variación entre grupos, y
	\item variación total, como combinación de las 
		dos anteriores.
	\end{enumerate}

	Consideraremos para realizar este test que disponemos de 
$n$ conjuntos de datos y $k$ clasificadores. Notaremos por 
$X_{ij}$ el rendimiento del algoritmo $j$ para la base de 
datos $i$. Así, $\bar{\bar{X}}$ se refiere a la media entre 
todos los valores de los clasificadores para las bases de
datos; $\bar{X}_{i \cdot}$ la media para la base de datos $i$ 
y $\bar{X}_{\cdot j}$ la media del clasificador $f_j$.\\
	El modelo para el rendimiento de un clasificador en una 
base de datos es
	\[
		X_{ij} = \bar{\bar{X}} + \alpha_i + e_{ij},
	\]
	donde $\alpha_i$ es la variabilidad del rendimiento entre 
los distintos conjuntos de datos, que representa una v.a. con media 0 y
varianza $\sigma_A^2$ y $e_{ij}$ representa la variabilidad
dentro del propio conjunto de datos, v.a. con media 
0 y varianza $\sigma^2$. Tanto $\alpha_i$ como
$e_{ij}$ se suponen normales. Además, se supone $e_{ij}$ 
independiente de $\alpha_i$ y de los demás $e_{ij}$.\\
	Definiremos en primer lugar algunas medidas de la 
variación, como la existente debida a los clasificadores
	\[ 
		SS_C = \sum\limits_{j=1}^k
				 \sum\limits_{i=1}^m
					\left( \bar{X}_{\cdot j} - 
						   \bar{\bar{X}} \right)^2 =
				n \sum\limits_{j=1}^k
					\left( \bar{X}_{\cdot j} - 
					  	   \bar{\bar{X}} \right)^2,
	\]
	de igual manera, la variación relativa a los conjuntos de
datos:
	\[ 
		SS_B = k \sum\limits_{i=1}^n
					\left( \bar{X}_{i \cdot} - 
					  	   \bar{\bar{X}} \right)^2.
	\]
	
	La variación total entre clasificadores se obtiene con
	\[
		SS_{Total} =  
			\sum\limits_{j=1}^k
				\sum\limits_{i=1}^m
					\left( \bar{X}_{ij} - 
						   \bar{\bar{X}} \right)^2.
	\]
			
	La variación en el error se deriva de la diferencia entre 
la variación total y la combinación de $SS_B$ y $SS_C$:
	\[
		SS_{Error} = SS_{Total} - (SS_B + SS_C).
	\]
	
	Definimos ahora los estadísticos a partir de los cuales 
realizaremos el test. En primer lugar, la variabilidad entre
clasificadores:
	\[
		MS_C = \frac{SS_C}{k-1},
	\]
	con $k-1$ el grado de libertad de $SS_C$. El segundo 
estadístico hace referencia a la variabilidad dentro de los 
clasificadores:
	\[
		MS_{Error} = \frac{SS_{Error}}{(n-1)(k-1)},
	\]
	con $(n-1)(k-1)$ los grados de libertad de $SS_{Error}$.
Finalmente, el estadístico con el que se lleva a cabo el
test, el cociente $F$ se obtiene como
	\[
		F = \frac{MS_C}{MS_{Error}},
	\]
	que sigue una distribución $F$ de Snedecor con $k-1$ y
$(n-1)(k-1)$ grados de libertad.

\paragraph{ANOVA de una vía} Se trata de un modelo más 
simple que no distingue entre la variabilidad dentro de los
conjuntos de datos y entre conjuntos de datos, sino que ambas
variabilidades se analizan de forma conjunta. En términos del
caso anterior, $X_{ij} = \bar{\bar{X}} + \alpha_i + e_{ij}$, 
con $\alpha_i = \bar{X}_{\cdot j} - \bar{\bar{X}}$ y $e_{ij}$
hace referencia al error aleatorio, que sigue una 
distribución normal, con media $0$ y varianza $\sigma^2$. \\
	Al calcular los estadísticos necesarios, tenemos que los 
grados de libertad de $SS_{Error} = SS_{Total} - SS_{C}$ son 
$nk-k$ y por consiguiente $MS_{Error} = \frac{SS_{Error}}
{nk-k}$. Este test asume la independencia entre las muestras
a diferencia del test ANOVA previo, donde la correlación de 
las muestras, esto es, el rendimiento de cada algoritmo para
el mismo conjunto de datos se tenía en cuenta. Además, este
test tiene una menor potencia.
	

\chapter{Test no paramétricos}
\label{chapter:test_no_parametricos}

	Para poder comprender este concepto, definimos la familia 
no paramétrica de distribuciones en términos de Pesarin 
(\cite{pesarin2010permutation}), sobre 
la que se realizan los test no paramétricos. 
 
\begin{definicion}[Familia de distribuciones no paramétrica]
	Una familia de distribuciones $\mathcal{P}$ se dice que 
se comporta de manera no paramétrica cuando no es posible
encontrar un espacio de dimensión finita $\Theta$ tal que hay 
una relación uno a uno entre $\Theta$ y $\mathcal{P}$, en el 
sentido en que cada elemento $P \in \mathcal{P}$ no puede ser 
identificado por un único elemento $\theta \in \Theta$ y 
viceversa.
\end{definicion}

	Si existiese esta relación, $\mathcal{P}$ sería una 
familia paramétrica y $\theta$ sería el parámetro.\\
	En base a esta definición, consideraremos la inferencia 
estadística no paramétrica aquella que se realiza sobre una 
muestra suponiendo la pertenencia de ésta a una familia no 
paramétrica de distribuciones. Un concepto relacionado es la 
propiedad de ser libre de distribución, del que tomamos la definición de Hollander (\cite{Hollander99}):
	
\begin{definicion}[Libre de distribución]
	Decimos que un estadístico es libre de distribución si 
para $H_0$ cierta, la distribución del estadístico no depende 
de la distribución $F$ de la población.
\end{definicion}
	
\section{Comparación con test paramétricos}
	
	En la inferencia clásica se efectúan suposiciones sobre 
la población de la que se extraen muestras para realizar la 
inferencia. Aunque estas suposiciones están normalmente 
justificadas, en ocasiones no se dan las circunstancias 
necesarias para aplicar estas técnicas o su uso no está bien 
documentado. Por ello surgen las técnicas no paramétricas. 
Nótese que es distinta la suposición al aplicar estas 
técnicas no paramétricas, pues consideramos en este caso que 
no conocemos su distribución.\\
 	La principal ventaja de los test no paramétricos es que  
las hipótesis son más generales con lo que se pueden aplicar 
en un mayor número de problemas. Una de las condiciones 
habituales es la continuidad, aunque hay otras condiciones 
más estrictas como que la población sea simétrica para según 
qué test. Esto repercute en que se le pueden aplicar 
funciones a las muestras obtenidas para la realización de los 
test, a diferencia de en los test paramétricos dado que las 
muestras deben generalmente provenir de una población de 
forma conocida.\\
	Hay test de hipótesis que no están relacionados con 
valores de parámetros (a diferencia de los paramétricos). Son 
más simples de aplicar, las matemáticas, menos sofisticadas y 
basadas en la combinatoria, están relacionadas con las 
propiedades usadas en el proceso inductivo. 
Además, las distribuciones asintóticas son 
distribuciones conocidas como la normal o la chi cuadrado. Al 
relajar las condiciones sobre los datos de entrada, es menos 
sensible al \textit{dirty-data}, datos con errores usados en 
el entrenamiento del clasificador. Esto implica una mayor 
robustez en los test no paramétricos.\\
 	Correctamente aplicados, los test paramétricos, al 
disponer de mayor información tienen una mayor potencia, sin 
embargo, cuando se disponen de menos datos, y por tanto es 
más difícil que se den las condiciones de los test 
paramétricos, la potencia es similar. Además, en esta 
situación, la falta de eficiencia de los test no paramétricos 
se suple con la falta de precisión de los test paramétricos 
al aproximar la distribución de los datos por la distribución 
asintótica conocida.\\


 	Se incluye en esta sección el desarrollo de los test no
paramétricos más utilizados en aprendizaje automático 
utilizando como referencia las indicaciones de Gibbons y 
Chakraborti (\cite{DBLP:reference/stat/GibbonsC11}) y Sheskin 
(\cite{sheskin2003handbook}).
 	
 	
\section{Test de aleatoriedad}

	Una de las condiciones para la realización de los test 
estadísticos, tanto de los paramétricos como de los no 
paramétricos, es la aleatoriedad de la muestra de partida. La 
hipótesis nula para los test que serán presentados en esta 
sección será la aleatoriedad de la muestra, mientras que la 
hipótesis alternativa será la presencia de un patrón. 
En el problema que nos ocupa, podríamos usar estos test 
en el estudio del efecto de un parámetro, partiendo de la
hipótesis nula de que dicho parámetro no afecta (y por tanto
los resultados son aleatorios) mientras que la hipótesis
alternativa contempla la existencia de un patrón conforme
este parámetro crece o decrece.

\subsection{Test basado en el número de rachas}

	Supongamos una secuencia de $n$ elementos de dos tipos, 
$n_1$ del primer tipo y $n_2$ del segundo, $n = n_1 + n_2$. 
Sea $R_1$ el número de rachas del primer tipo, $R_2$ el 
número de rachas del segundo tipo, $R = R_1 + R_2$. Siendo 
cierta la hipótesis nula (la aleatoriedad de la muestra), 
procedemos a obtener la distribución de $R$.
	
\begin{lema} 
	El número de formas distintas de distribuir $n$ objetos 
en $r$ posiciones consecutivas es ${n-1 \choose r-1}, n \geq 
r, r \geq 1$.
\end{lema}

\begin{teorema}
	Sean $R_1$ y $R_2$ los números de rachas de los $n_1$ de 
tipo 1 y los $n_2$ elementos de tipo 2 respectivamente en una 
muestra de tamaño $n = n_1 + n_2$. La función de distribución 
de probabilidad conjunta de $R_1$ y $R_2$ es
	\[ f_{R_1,R_2} (r_1, r_2) = 
		\frac{c {n_1 - 1 \choose r_1 - 1} 
				{n_2 - 1 \choose r_2 - 1}}
			{{n_1 + n_2 \choose n_1}}\;
		\begin{array}{l}
			r_1 = 1,2, \dots, n_1 \\
			r_2 = 1,2, \dots, n_2 \\
			r_1 = r_2 \text{ ó } r_1 = r_2 \pm 1
		\end{array}
	\]
	donde $c=2$ si $r_1 = r_2$ (hay igual número de rachas de 
elementos del tipo 1 y del tipo 2) y $c=1$ si $r_1 = 
r_2 \pm 1$ (hay una racha más del tipo 1 ó 2).
\end{teorema}

	Para muestras de un mayor tamaño (aquellas en el que 
$n_1, n_2 \geq 10$) se suele utilizar una aproximación 
utilizando la distribución asintótica supuesto cierta 
$H_0$.\\
	Suponemos que el tamaño de la muestra $n \rightarrow 
\infty$, de forma en que $\frac{n_1}{n} \rightarrow \lambda$, 
$0<\lambda<1$. De aquí obtenemos
	\[ \underset{n \rightarrow \infty}{\lim} \mathbb{E}[R/n] = 
			2\lambda (1-\lambda) 
				\underset{n \rightarrow \infty}{\lim} 
					\var(R\sqrt{n}) =
			4\lambda^2(1-\lambda)^2
	\]
	
\begin{teorema}
	La distribución de probabilidad de $R$, es decir, el 
número total de rachas en una muestra aleatoria es:
	
	\begin{equation}
		f_R(r) = \left\lbrace\begin{array}{ll}
	2 {n_1-1 \choose r/2-1} {n_2-1 \choose r/2-1} 
		\big/ {n_1 + n_2 \choose n_1} &
			\textit{ si } r \text{ es par} \\
	\left[
		{n_1-1 \choose (r-1)/2} {n_2-1 \choose (r-3)/2} +  
		{n_1-1 \choose (r-3)/2} {n_2-1 \choose (r-1)/2} 
	\right]
		\big/ {n_1 + n_2 \choose n_1} &
			\textit{ si } r \text{ es impar} \\		
		\end{array}\right.
	\label{th-dist-R}
	\end{equation}
	para $r=2, 3, \dots, n_1 + n_2.$
\end{teorema}
	
	Si llamamos $Z = \frac{R - 2n\lambda (1-\lambda)}
{2 \sqrt{n}\lambda (1-\lambda)}$ y sustituimos en 
\ref{th-dist-R}, obtenemos la distribución estandarizada de 
$R$, $f_Z(z)$. Entonces aplicamos la fórmula de Stirling y el 
límite queda de la forma
	\[ \underset{n \rightarrow \infty}{\lim} \log f_Z(z)=
			-\log \sqrt{2\pi} - \frac{1}{2} z^2,	\]
	con lo que la distribución límite de $Z$ es la normal. Por
ejemplo, si $z_{\alpha/2}$ es el valor límite de la distribución
normal siendo $\alpha$ el nivel de significación, rechazaremos
la hipótesis nula si $\abs{Z} \geq z_{\alpha/2}$.
	
	
\subsubsection{Test basado en rachas crecientes y decrecientes}	

	Para este test consideramos una serie de datos de tipo 
numérico ordenados temporalmente y queremos comprobar la 
hipótesis de la aleatoriedad de la muestra.\\
	Para una muestra de $n$ elementos, supongamos que podemos 
ordenarlos de la forma $a_1 < \dots < a_n$ (estamos 
suponiendo que no hay dos iguales). Si la hipótesis nula fuese 
cierta, nuestra muestra se corresponderá con una de las $n!$ 
permutaciones con igual probabilidad. Usaremos para este test 
las rachas crecientes y decrecientes. Construimos la 
secuencia $D_{n-1}$, cuyo elemento $i$-ésimo es el signo de 
$x_{i+1} - x_i,\ i=1, \dots, n-1$. Sean $R_1, \dots, R_{n-1}$ 
el número de rachas de longitud $1, \dots, n-1$ 
respectivamente. $f_n(r_{n-1}, \dots, r_1)$ indica la 
probabilidad de obtener $r_j$ rachas de longitud $j$ supuesta 
cierta la hipótesis nula. Escribiremos como $u_n$ la 
frecuencia absoluta $f_n = \frac{u_n}{n!}$. Para obtener la 
función de distribución, partiremos del caso $n=3$.\\
	Sean $a_1 < a_2 < a_3$. La distribución de probabilidad 
será 
	\[ 
	f_3(r_2, r_1) = 
		\left\lbrace\begin{array}{cc}
			\frac{2}{6} & \text{si } r_2 = 1, r_1 = 0 \\
			\frac{4}{6} & \text{si } r_2 = 0, r_1 = 2 
		\end{array}\right.
	\]
	dado que las únicas rachas de longitud 2 son 
$(a_1, a_2, a_3) \rightarrow (+,+)$ y $(a_3, a_2, a_1)$ 
$\rightarrow (-,-)$, siendo las demás posibles combinaciones 
dos rachas de longitud 1.\\
	Las posibilidades a la hora de insertar un elemento $a_n$ 
en las permutaciones de $S_n$ son:
\begin{enumerate}
	\item Se añade una racha de longitud 1.
	\item Una racha de longitud $i-1$ se convierte en una 
		de longitud $i$, $i=2,\dots n-1$.
	\item Una racha de longitud $h=2i$ se convierte en una de 
		longitud $i$, seguida por otra de longitud $1$, 
		seguida por otra de longitud $i$.
	\item Una racha de longitud $h=i+j$ se convierte en
	\begin{enumerate}
		\item Una racha de longitud $i$ seguida por otra de 	
			longitud $1$ seguida por otra de longitud $j$,
		\item Una racha de longitud $j$ seguida por otra de 	
			longitud $1$ seguida por otra de longitud $i$,
	\end{enumerate}
		con $h>i>j$, $3 \leq h \leq n-2$.
\end{enumerate}

	De forma general, la frecuencia $u_n$ conocido $u_{n-1}$ 
sigue la siguiente relación:
\begin{align*}
	& u_n (r_{n-1}, \dots, r_h, \dots, r_i, \dots, r_j, 
			\dots, r_1)= 
		2 u_{n-1}(r_{n-2}, \dots, r_1-1) \\
	&+ \sum\limits_{i=2}^{n-1} 
		(r_{i-1} + 1)
		u_{n-1}(r_{n-2},\dots, r_i-1, r_{i-1}+1,\dots, r_1)\\
	&+ \sum\limits_{i=1, h=2i}^{\lfloor (n-2)/2 \rfloor} 
		(r_{h} + 1)
		u_{n-1}(r_{n-2},\dots, r_h+1,\dots r_i-2,
				\dots, r_1-1)\\
	&+ 2 \underbrace{\sum\limits_{i=2}^{n-3} 
		\sum\limits_{j=1}^{i-1}}_{h=i+j, h \leq n-2}
		(r_{h} + 1)
		u_{n-1}(r_{n-2},\dots, r_h+1,\dots, r_i-1,
				\dots, r_1-1)			
\end{align*}
	
	Otro test que se podría realizar es el del número total 
de rachas, independientemente de su longitud. El número de 
rachas total sería $R = \sum\limits_{i=1}^{n-1} R_i$. Usando 
el procedimiento anterior, se llega a que la distribución 
del estadístico $\frac{R-\mu}{\sigma}$ con $\mu = 
\frac{2n-1}{3}$ y $\sigma^2=\frac{16n-29}{90}$ es la 
normal estándar.
	
\section{Test de bondad del ajuste}

	Una cuestión relevante en la estadística es la obtención 
de la forma de la población de la que se obtiene una muestra. 
En los test paramétricos, como se ha mencionado previamente, 
es habitual que la información relativa a la forma de la 
población se incluya entre las condiciones, por ejemplo en el 
test basado en la distribución $t$ de Student es necesaria la 
normalidad de la población. Por consiguiente, nos interesa 
disponer de test que nos indiquen cómo de confiable es la 
hipótesis de la normalidad para nuestra muestra. Los test de 
bondad del ajuste se ocupan de la forma de la población y no 
de la distribución concreta incluyendo parámetros de escala o 
localización. La utilización de estos test en nuestro problema
nos permiten, según los resultados, justificar el
uso de test paramétricos para realizar la comparación. Tras
la aplicación de estos test podríamos, si los datos
fuesen consistentes con la hipótesis de normalidad,
decidirnos por la aplicación del test ANOVA en lugar
de uno no paramétrico.
	
\subsection{Test chi cuadrado}

	Sea una muestra aleatoria de tamaño $n$ obtenida de una 
población con función de distribución desconocida $F_X$. El 
test tendrá como hipótesis nula
		\[ H_0: F_X(x) = F_0(x) \ \forall x \]
	con $F_0$ conocida contra
		\[ H_1: F_X(x) \neq F_0(x) \text{ para algún }  x \]
	Para realizar este test, los datos deben disponerse en
categorías, bien mediante los valores que toma la variable 
para distribuciones discretas o mediante rangos especificados 
al realizar el experimento para distribuciones continuas. Una 
vez realizadas estas categorías podemos obtener las 
frecuencia esperada si la hipótesis nula es cierta de 
$F_0(x)$.\\
	
	Supongamos que disponemos de $n$ muestras clasificadas en 
$k$ categorías mutuamente excluyentes. Sea $f_i$ la 
frecuencia observada y $e_i$ la frecuencia esperada para cada 
muestra. El estadístico que definimos para la realización de 
este test es $Q = \sum\limits_{i=1}^k 
\frac{(f_i-e_i)^2}{e_i}$, para el cual estudiaremos su 
distribución asintótica.\\
	
	Sean $\theta_1, \dots, \theta_k$ las probabilidades de 
pertenencia a cada clase y $f_1, \dots, f_k$ los valores 
observados. La función de verosimilitud será:
	
	\[ L(\theta_1, \dots, \theta_k) = 
			\prod\limits_{i=1}^k \theta_i^{f_i},\			
	   \text{ con } f_i = 0, 1, \dots, n; \
	   \sum\limits_{i=1}^k \theta_i = 1, \
	   \sum\limits_{i=1}^k f_i = n
	 \]
	 
	Podemos por tanto escribir la hipótesis nula de la 
siguiente forma:
	 
	\[ H_0 = \theta_i = \theta_i^0,\ i = 1, \dots, k \]
	 
	habiendo obtenido cada $\theta_i^0$ de $F_0$. El 
estimador de máxima verosimilitud es $\hat{\theta}_i = 
\frac{f_i}{n}$. Entonces la razón de verosimilitud es
	 
	 \[ 
	 T = \frac{L(\hat{\omega})}{L(\hat{\Omega})}
	   = \frac{L(\theta_1^0, \dots, \theta_k^0)}
	   		{L(\hat{\theta}_1^0, \dots, \hat{\theta}_k^0)}
	   = \prod\limits_{i=1}^k
	   		\left( 
	 			\frac{\theta_i^0}{\hat{\theta_i}} 
	 		\right)^{f_i}
	 \]
	 y la distribución de la v.a. $-2 \log T$ puede ser 
aproximada por la distribución chi cuadrado. Debido a que 
estimamos $k-1$ parámetros (el parámetro restante se deduce 
de la restricción $\sum\limits_{i=1}^k \theta_i = 1$). 
Tenemos entonces
	 \begin{equation}
	 2 \log T = 
	 		-2 \sum\limits_{i=1}^k
	 			f_i \left(
	 					\log \theta_i^0 - \log \frac{f_i}{n}
	 				\right)
	 \label{2logT}
	 \end{equation}
	y mostramos a continuación que esta expresión es 
equivalente asintóticamente a la expresión de $Q$.\\
	
	La serie de Taylor de $\log \theta_i$ centrada en 
$\hat{\theta}_i = f_i/n$ es
	\[ \log \theta_i = 
			\log \hat{\theta}_i +
			(\theta_i - \hat{\theta}_i)
				\frac{1}{\hat{\theta}_i} +
			\frac{(\theta_i - \hat{\theta}_i)^2}{2!}
				\left(-\frac{1}{\hat{\theta}_i^2}\right) +
			\dots
	\]
	
	con lo que
	
	\begin{align*}	
	 \log \theta_i^0 - \log \frac{f_i}{n} & = 
			\left(
				\theta_i^0 - \frac{f_i}{n}
			\right) \frac{n}{f_i} -
			\left(
				\theta_i^0 - \frac{f_i}{n}
			\right)^2 \frac{n^2}{2f_i^2} + \epsilon \\
		&= \frac{n\theta_i^0 - f_i}{f_i} -
			\frac{(n\theta_i^0 - f_i)^2}{2f_i^2} +\epsilon
	\end{align*}
	
	con $\epsilon = \sum\limits_{j=3}^\infty
			(-1)^{j+1} 
			\left( \theta_i^0 - \frac{f_i}{n}\right)^j
			\frac{n^j}{j!f_i^j}$. Sustituyendo en \ref{2logT} 
nos queda
			
	\begin{align*}
	-2 \log T &= 
		-2 \sum\limits_{i=1}^k n\theta_i^0 - f_i -
		\sum\limits_{i=1}^k 
			\frac{(n\theta_i^0 - f_i)^2}{f_i} + 
		\sum\limits_{i=1}^k \epsilon' = \\
	&= 0 + 
	   \sum\limits_{i=1}^k \frac{(f_i-e_i)^2}{f_i} +
	   \epsilon''	
	\end{align*}	 
	
	Por la ley de los grandes números, $F_i/n$, 
para $F_i$ la frecuencia de las clases, es un 
estimador consistente de $\theta_i$, es decir
	
	\[ \lim_{n \rightarrow \infty} \left[
			P \left(
				\abs{ \frac{F_i}{n} - \theta_i } 
					> \varepsilon
			\right) \right] = 0 \; \forall \varepsilon > 0 \]
	
	\subsubsection{Kolmogorov-Smirnov}
	
	En el test anterior sólo se realizaban $k$ comparaciones 
a pesar de disponer de $n$ ($n \geq k$) observaciones. Si las 
$n$ muestras provienen de una distribución continua, 
podríamos realizar la comparación entre la función de 
distribución que constituye la hipótesis nula y la 
\textbf{función de distribución empírica} (edf) $S_n$ 
definida de la siguiente forma:
	\[ S_n(x) = \frac{\text{número de muestras} \leq x}{n} \]
	Para una definición formal introducimos notación sobre 
estadísticos ordenados. Sea $X_1, \dots, X_n$ una muestra 
aleatoria de una población con función de distribución 
continua $F_X$. Sea $X_{(1)}$ el menor valor de $X_1, \dots, 
X_n; X_{(2)}$ el segundo menor;$ \dots; X_{(n)}$ el mayor. 
Entonces se puede definir $S_n(x)$ como:
	\[ S_n(x) = 
		\left\lbrace\begin{array}{ll}
			0 & \text{si } x < X_{(1)} \\
			i/n & \text{si } X_{(i)} < x <X_{(i+1)},
				i = 1, \dots, n-1 \\
			1 & \text{si } x > X_{(n)} \\
	\end{array}\right.
	\]
	
	$S_n(x)$ es un estimador consistente para $F_X(x)$ y 
conforme $n$ crece, se aproxima a $F_X(x)$ para todo $x$. 
Entonces cabe esperar que el error vaya disminuyendo. Se 
define el estadístico
	\[ D_n = \underset{x}{\sup} 
				\vert S_n(x) - F_0(x) \vert, \]
	que si la hipótesis nula es cierta es un buen indicador 
de la precisión de la estimación. El estadístico $D_n$ es 
especialmente útil en la inferencia no paramétrica debido a 
que la distribución de $D_n$ no depende de $F_0(x)$ mientras 
sea continua. Se definen las desviaciones direccionales
\begin{align*}
		D_n^+ &= \underset{x}{\sup} [S_n(x) - F_0(x)]  		\\
		D_n^- &= \underset{x}{\sup} [F_0(x) - S_n(x)]
\end{align*}	 
	
\begin{teorema}
	Los estadísticos $D_n, D_n^+$ y $D_n^-$ son libres de 
distribución para cualquier función de distribución continua
$F_0$.
\end{teorema}
	
\begin{teorema}
	Si $F_X$ es una función de distribución continua, 
entonces para cada $d>0$
	\[ \underset{n \rightarrow \infty}{lim}
			P(D_n \leq d/\sqrt{n}) = L(d) \]
	con
	\[ L(d) = 1 - 2 \sum\limits_{i=1}^\infty 
			(-1)^{i-1} e^{-2i^2d^2}	\]
\end{teorema}

\begin{teorema}
	Si $F_0$ es una función de distribución continua, 
entonces bajo $H_0$ para cada $d>0$
	\[ \underset{n \rightarrow \infty}{lim}
			P(D_n^+ < d/\sqrt{n}) = 1-e^{-2d^2} \]
\end{teorema}	
	Como consecuencia de este teorema podemos usar las tablas
de la distribución chi cuadrado:
	
\begin{corolario}
	Si $F_0$ es una función de distribución continua, 
entonces para cada $d \geq 0$, la distribución límite de 
$V = 4n {D_n^+}^2$ para $n \rightarrow \infty$ es la 
distribución chi cuadrado con dos grados de libertad.
\end{corolario}	
	
	Para aplicar el test, partimos de $\tau_\alpha$ el valor
crítico para el nivel de significación $\alpha$ de 
$\chi^2_2$, entonces del corolario anterior obtenemos el 
valor límite para $D^+_n$ despejando de $\tau_\alpha
= 4n {D_n^+}^2$. Si el $D^+_n$ de la muestra fuese mayor,
rechazaríamos la hipótesis nula.
	
\section{Análisis del conteo de datos}
	
	A la hora de comparar el rendimiento de algoritmos, una 
primer enfoque puede ser contar el número de casos en los que 
un algoritmo acierta y otro no. Para representar los datos 
necesarios para este tipo de test, se utiliza una matriz de 
contingencia $r \times k$, donde cada posición se corresponde 
con el número de muestras observadas que cumple las 
características de la fila y la columna.
	
\subsection{Test de McNemar}
	
	En este test suponemos una tabla $2 \times 2$ (Tabla \ref{Tb:McNemar}) que recoge 
la respuesta de $n$ elementos ante una cierta prueba antes y 
después de un tratamiento o, en el caso en el que estamos, si 
los $n$ individuos de la muestra están bien o mal 
clasificados por dos algoritmos de clasificación.
	
\begin{table}[H]
\centering
\caption{Tabla de contingencia de McNemar}
\label{Tb:McNemar}
\begin{tabular}{|lc|cc|}
\hline
                                                          &       & \multicolumn{2}{c|}{Clasificador $f_2$} \\ \cline{3-4} 
                                                          &         & Fallo              & Acierto            \\ \hline
\multicolumn{1}{|c|}{\multirow{2}{*}{Clasificador $f_1$}} & Fallo   & $c_{00}$           & $c_{01}$           \\
\multicolumn{1}{|c|}{}                                    & Acierto & $c_{10}$           & $c_{11}$           \\ \hline
\end{tabular}
\end{table}

	En este test la hipótesis nula es que la probabilidad de 
que un individuo sea bien clasificado es la misma para ambos 
algoritmos. Notamos por $\theta_{ij},\ i=0,1,\ j=0,1$ la 
probabilidad de pertenecer a cada grupo. Notamos 
$\theta_{1 \cdot} = \theta_{10} + \theta_{11}$, 
$\theta_{\cdot 1} = \theta_{01} + \theta_{11}$, la 
probabilidad de estar bien clasificado por el primer y el 
segundo clasificador respectivamente. La hipótesis nula se 
escribe entonces como $H_0: \theta_{1 \cdot} = 
\theta_{\cdot 1}$, para la cual se considera el estadístico 
basado en $T = \frac{c_{1 \cdot} - c_{\cdot 1}}{N}$ como 
estimador insesgado. Podemos escribir $T = \frac{c_{01} - 
c_{10}}{N}$. El test de McNemar está basado en
\begin{equation}
\label{eq:McNemarChi}
	\chi = \frac{(c_{01} - c_{10})^2}{(c_{01} + c_{10})}
\end{equation} 
	que sigue aproximadamente una distribución $\chi^2$ con 
un grado de libertad. Obtengamos ahora la esperanza y 
varianza de $T$. Puesto que $c_{01}$ y $c_{10}$ siguen una 
distribución binomial con parámetros $n$, $\theta_{01}$ y 
$\theta_{10}$ respectivamente, obtenemos

\begin{align*}
	\mathbb{E}[c_{01}] 	&= n \theta_{01} \\
	\var(c_{01}) &= n \theta_{01} (1 - \theta_{01}) \\
	\mathbb{E}[c_{10}] 	&= n \theta_{10} \\
	\var(c_{10}) &= n \theta_{10} (1 - \theta_{10}).
\end{align*}

	La varianza de $T$ se puede ver obtener de 
	
\begin{equation}
	N^2 \var(T) = \var(c_{01}) + \var(c_{10}) - 2 cov(X_{01}, c_{10}).
	\label{eq:McNemarVar}
\end{equation}

	Para el término de la covarianza podemos usar que la 
distribución de $(c_{00}, c_{01}, c_{10}, c_{11})$ es una 
distribución multinomial con probabilidades $(\theta_{00}, 
\theta_{01}, \theta_{10}, \theta_{11})$. Entonces, haciendo 
uso de la función generatriz de momentos de $c_{01}$ y 
$c_{10}$,
	\[ 
		M_X(t_1, t_2) = 
		\left(
			\theta_{01} e^{t_1} + 
			\theta_{10} e^{t_2} + 
			[ 1 - (\theta_{01} + \theta_{10})] \right)^n,
	\]
obtenemos $\mathbb{E}[c_{01}c_{10}] = \frac{\partial^2 }
{\partial t_1 \partial t_2} M_X(0,0) = n(n-1)\theta_{01} 
\theta_{10}$ y por tanto escribimos la covarianza como:
	\[
		\cov( c_{01}, c_{10} ) =
			n (n-1) \theta_{01} \theta_{10} -
			(n \theta_{01}) (n \theta_{10}) = 
			- n \theta_{01} \theta_{10}.
	\]
	
	Sustituyendo en \ref{eq:McNemarVar}, tenemos que $T$ 
tiene media $\theta_{01}-\theta_{10}$ y $\var(T) = 
\frac{(\theta_{01} + \theta_{10}) - (\theta_{01} - 
\theta_{10})^2 }{n}$. Bajo la hipótesis nula $\theta_{01} = 
\theta_{10}$, $\mathbb{E}[T] = 0$ y $\var(T) = \frac{\theta_{01}+
\theta_{10}}{n}$, que se puede estimar de manera consistente 
a partir del estadístico $(c_{01} + c_{10})/n^2$.\\
	Otro enfoque es observar el número de parejas 
discordantes $S = c_{01} + c_{10}$. Entonces $c_{01}$ sigue 
una binomial con parámetros $S$ y $\frac{\theta_{01}}
{\theta_{01}+\theta_{10}}$. Es posible que estemos 
interesados en una alternativa unilateral, esto es $H_1: 
\theta_{\cdot 1} > \theta_{1 \cdot}$, en la que la hipótesis 
alternativa supone que un algoritmo de clasificación es mejor 
que otro. Para esta alternativa el $p$-valor exacto viene 
dado por $\sum\limits_{j=0}^S {S \choose j}(0.5)^S$. Para 
muestras de un gran tamaño, podemos obtener un $p$-valor 
aproximado basándonos en $Z = \frac{c_{01} - c_{10}}
{\sqrt{c_{01} + c_{01}}}$, que sigue aproximadamente una 
distribución normal estándar.
	
\section{Test basados en una muestra y en muestras apareadas}
	
	En este apartado presentaremos los test no paramétricos 
análogos a test paramétricos como el test $t$ de Student para 
las hipótesis nulas $H_0: \mu = \mu_0$ para una única muestra 
y $H_0: \mu_X - \mu_Y = \mu_0$ para muestras apareadas. En 
los test no paramétricos sólo serán necesarias condiciones 
sobre la continuidad de la población. 
	
\subsection{Test de signo}

	Supongamos una muestra de $n$ elementos de una población 
con mediana desconocida $M$, donde suponemos que $F_X$ 
es continua y estrictamente creciente al menos en el entorno 
de $M$. Esto significa que $F_X^{-1}(0.5) = M$. La hipótesis 
nula a comprobar se corresponde con el valor de la mediana:

	\[ H_0: M = M_0	\text{ o equivalentemente }
			P(X > M_0) = 0.5 \]
	con $M_0$ un valor dado. Como hemos supuesto que $F_X$ 
tiene una única mediana, la hipótesis significa que $M_0$ 
divide el área de la función de densidad en dos partes 
iguales. Denotamos por $K$ el número de observaciones mayores 
que $M_0$. Podemos considerar entonces que estamos obteniendo 
una muestra de una v.a. $K$ que sigue una distribución de 
Bernouilli con parámetros $n$ y $\theta=P(X>M_0)$, y 
$\theta=0.5$ si la hipótesis nula es cierta. Se llama test de 
signo debido a que $K$ es el número de signos positivos en 
$X_i - M_0, i = 1, \dots, n$. La hipótesis alternativa queda 
por tanto
	\[ H_1: M \neq M_0	\text{ ó }
			\theta = P(X > M_0) \neq 0.5 \]
	La región crítica, para $\alpha$ el nivel de 
significación, es $K \geq k_{\alpha/2}$ ó $K \leq 
k_{\alpha/2}'$, con $k_{\alpha/2}$ y $k_{\alpha/2}'$ el menor 
y el mayor entero respectivamente tales que
	\[ \sum\limits_{i=k_{alpha/2}}^n
			{n \choose i}(0.5)^n \leq \frac{\alpha}{2}
		\text{ y }
		\sum\limits_{i=0}^{k_{alpha/2}'}
			{n \choose i}(0.5)^n \leq \frac{\alpha}{2}
	\]

\subsubsection{Test de rangos con signos de Wilcoxon}

	El test anterior usa únicamente el signo de la diferencia 
de la muestra a la mediana, sin considerar la distancia. Para 
este test sí consideramos la distancia a la mediana aunque 
necesitamos suponer la simetría de la población. \\
	Sea $X_1, \dots, X_n$ una muestra de una función de 
distribución continua $F$ con mediana $M$. De ser cierta la 
hipótesis nula $H_0: M = M_0$ las diferencias $D_i = X_i - 
M_0$ estarían distribuidas de manera simétrica en torno 
a 0.\\
	Supongamos que ordenamos las diferencias absolutas 
$|D_1|, \dots, |D_n|$ del valor absoluto más pequeño al más 
grande y le asignamos los puestos $1, 2, \dots, n$. Sea $T^+$ 
el valor esperado de la suma de los puestos con diferencias 
positivas, $T^-$ la de aquellos con diferencias negativas. 
Supuesta cierta la hipótesis nula, al ser la población 
simétrica, $T^+ = T^-$. Para realizar el test, notaremos por
$r(\cdot)$ la función que asigna el puesto de la v.a.. 
Definimos
	\[ T^+ = \sum\limits_{i=1}^n Z_i r(|D_i|); \quad
	   T^- = \sum\limits_{i=1}^n (1-Z_i) r(|D_i|) \]
	   
	 con $ Z_i = \mathbb{I}[ D_i > 0 ]$. Entonces,
	 		
	\[ T^+  - T^- = 2 \sum\limits_{i=1}^n
					 Z_i r(|D_i|) - \frac{N(N+1)}{2} \]
	
	Bajo la hipótesis nula, $Z_i$ sigue una distribución de 
Bernouilli con $\mathbb{E}(Z_i) = 1/2$ y $\var(Z_i) = 1/4$. Usando que 
$T^+$ es una combinación lineal de las variables $Z_i$, 
tenemos
	\[ \mathbb{E}[ T^+ | H_0 ] = \sum\limits_{i=1}^n 
						\frac{r(|D_i|)}{2} 
					= \frac{n(n+1)}{4}			\]
	y
	
	\[ \var( T^+ | H_0 ) = \sum\limits_{i=1}^n 
						\frac{[r(|D_i|)]^2}{4} 
					= \frac{n(n+1)(2n+1)}{24}	\]
					
	Para la aplicación de este test se obtiene el estadístico 
$T = \min \{ T^+, T^- \}$. En la tabla de los valores 
críticos de $T$ para el test de rangos con signos de Wilcoxon 
(por ejemplo en \cite[Tabla A5]{sheskin2003handbook}) se encuentran aquellos valores 
para nivel de significación 0.05 y 0.01 para los cuales 
debemos rechazar la hipótesis nula si $T$ es menor que el 
valor correspondiente en la tabla.\\
	
\paragraph{Valores iguales a la mediana} En este test se han 
considerado los valores iguales a la mediana como negativos. 
Como las diferencias con la mediana están ordenadas de manera 
creciente, estos valores estarán necesariamente al principio 
y el impacto será menor. Sin embargo, podemos considerar la 
siguiente modificación,

	\[ Z_i = \left\lbrace \begin{array}{ll}
	 				1 & \text{si } D_i > 0 \\
	 				0 & \text{si } D_i < 0 \\
	 				1/2 & \text{si } D_i = 0 
	 		\end{array} \right. 					\]
	
	para repartir así entre $T^+$ y $T^-$ estas puntuaciones. 
	
\subsubsection{Tratamiento de empates}
	
	Al realizar la suposición de que la muestra se obtiene de 
una población continua, la probabilidad teórica de obtener 
dos valores idénticos es nula. Sin embargo en la práctica 
podemos obtener empates debido a que la población sea 
discreta o limitaciones en la precisión. Presentamos a 
continuación algunos métodos para solventar esta situación:
	
\paragraph{Aleatorización} Se selecciona un orden aleatorio 
para los elementos con igual valor.

\paragraph{\textit{Midranks}} Asigna a cada individuo de un 
grupo de valores empatados la media de los puestos de la 
clasificación que tendrían de ser distintos. Es decir, si hay 
tres valores con la tercera menor distancia a la mediana, 
ocuparían los puestos 3,4 y 5 y por tanto el valor asignado a 
cada uno de ellos sería $\frac{3+4+5}{3}=4$. Es el método más 
utilizado debido a su simplicidad. 

\paragraph{Estadístico medio} Se realiza el test estadístico 
para todas las posibles asignaciones para los términos 
empatados y se realiza la media de estos test.

\paragraph{Estadístico menos favorable} Habiendo encontrado 
todos los posibles valores del test, se escoge aquel con 
menor probabilidad de rechazar la hipótesis nula. Es la 
opción más conservadora al minimizar la probabilidad de 
cometer un error de tipo I.

\paragraph{Rango de probabilidad} Se devuelve el valor menos 
favorable a rechazar la hipótesis nula y el más favorable. No 
conduce a una única decisión a no ser que ambos valores 
caigan en la región crítica o fuera de ella.

\paragraph{Omisión de los valores empatados} Otra posibilidad 
es descartar los valores empatados. Conlleva una pérdida de 
información y generalmente introduce un sesgo hacia rechazar 
la hipótesis nula.

	
\section{Medidas de asociación en clasificaciones múltiples}	

	En este apartado presentamos la versión no paramétrica 
para el problema del análisis de la varianza. Enfocado al 
problema abordado en este trabajo, nos será muy útil para 
comparar el rendimiento de varios algoritmos en distintas 
conjuntos de datos. Los datos se presentan en una tabla $n \times 
k$ con entradas $X_{ij}$. En esta distribución influyen dos 
factores que llamaremos el efecto por filas y por columnas.\\
	Para nuestro problema, consideraremos que las filas se 
corresponden con conjuntos de datos y que las columnas 
constituyen los algoritmos utilizados. 
	
\subsection{Análisis bidimensional de la varianza mediante clasificaciones en una tabla $k \times n$ y comparaciones múltiples de Friedman}

	Partimos de una matriz $n \times k$. Consideraremos que 
las filas son independientes, pero no las columnas (en 
nuestro problema esto significa que las conjuntos de datos sí 
son independientes pero en cada algoritmo los resultados guardan 
correlación). La hipótesis nula que consideraremos será la 
equivalencia del rendimiento de todos los algoritmos.
Friedman sugirió cambiar cada valor de la 
matriz por la clasificación de cada fila. Notaremos por $R_{ij}$
el orden del algoritmo $j$ para el problema $i$.

	\[ \left( \begin{matrix}
		R_{11} & R_{12} & \dots & R_{1k} \\
		\vdots & \vdots & \ddots & \vdots \\
		R_{n1} & R_{k2} & \dots & R_{nk}
		\end{matrix} \right)	\]

	Entonces, $R_{i1}, \dots, R_{ik}$ es una permutación de 
los $k$ primeros números (salvo en caso de empate). 
Escribimos como $R_j$ el total para la columna $j$. 
Consideraremos la distribución de la muestra de la v.a. 
	\begin{equation}
		 S = \sum\limits_{j=1}^k
				\left[
					R_j - \frac{n(k+1)}{2}
				\right]^2 =
			\sum\limits_{j=1}^k \left[
				\sum\limits_{i=1}^n \left(
					R_{ij} - \frac{k+1}{2}
				\right)
			\right]^2,
	\label{S-Friedman}
	\end{equation}
bajo la hipótesis nula $H_0: \theta_1 = \dots = 
\theta_k$. En el caso de que no hubiera empates por filas 
habría un total de $(k!)^n$ posibles entradas igualmente 
probables. Las posibilidades se pueden enumerar y calcular el 
valor de $S$ para cada una de ellas, con lo que la 
distribución de probabilidad de $S$ sería $f_S(s) = 
\frac{u_s}{(k!)^n}$, con $u_s$ el número de asignaciones que 
tienen como suma de los cuadrados de las desviaciones totales 
por columnas $s$. Existe un método para obtener $u_s$ de $k$ 
y $n$ a partir de $k$ y $n-1$, y tablas con valores. Para los 
valores que exceden estas tablas se usa una aproximación de 
la distribución.\\
	Notando $\mu = \frac{k+1}{2}$, escribimos 
	\ref{S-Friedman}:
	
	\begin{align}
		S 	&= 	\sum\limits_{j=1}^k
					\sum\limits_{i=1}^n (R_{ij}-\mu)^2
				+ 2 \sum\limits_{j=1}^k
					\sum\limits_{i=1}^{n-1}
						\sum\limits_{p=i+1}^n
							(R_{ij}-\mu)(R_{pj}-\mu) 
				\nonumber \\
			&=	k \sum\limits_{j=1}^n (j-\mu)^2 + 2U 
				\nonumber \\
			&=  \frac{nk(k^2-1)}{12} + 2U	.
	\label{S-Friedman2}
	\end{align}
	
	Los momentos de $S$ vienen determinados por los momentos 
de $U$, que se derivan de:
	\[ 	\mathbb{E}[ R_{ij} ]= \frac{k+1}{2};\ 
		\var(R_{ij})=\frac{k^2-1}{12};\
		\cov(R_{ij},R_{iq}) = -\frac{k^2-1}{12}
	\]
	
	Además, dado que las observaciones en las diferentes 
filas son independientes, para $i \neq p$, se tiene
$\cov(R_{ij},R_{pq})=0$. Entonces
	\[ \mathbb{E}(U) = k {n \choose 2} \cov(R_{ij},R_{pq})=0, \] 
	con lo que $\var(U) = k\mathbb{E}(U^2)$, con
	\begin{align*}
		U^2 =& \sum\limits_{j=1}^k
				\sum\limits_{1 \leq i < p \leq n}
					(R_{ij} - \mu)^2 (R_{pj} - \mu)^2 \\
			+& 2 
			    \sum\limits_{j=1}^{k-1}
			      \sum\limits_{q=j+1}^k
					\sum\limits_{i=1}^{n-1}
					  \sum\limits_{p=i+1}^n	
					  	\sum\limits_{r=1}^{n-1}
						  \sum\limits_{s=r+1}^n
				(R_{ij} - \mu)(R_{pj} - \mu)
				(R_{rq} - \mu)(R_{sq} - \mu).
	\end{align*}
	
	Como $R_{ij}$ y $R_{pq}$ son independientes para $i \neq p$, tenemos
	\begin{align*}
		\mathbb{E}[U^2] =& \sum\limits_{j=1}^k
					\sum\limits_{1 \leq i < p \leq n}
						\var(R_{ij})\var(R_{pj}) \\
			 &+ 2 
			    \sum\limits_{j=1}^{k-1}
			      \sum\limits_{q=j+1}^n
			    {n \choose 2}
				\cov(R_{ij}, R_{iq}) \cov(R_{pj}, R_{pq}) \\
			=& k {n \choose 2} \frac{(k^2-1)^2}{144} +
				2{k \choose 2}{n \choose 2} 
					\frac{(k^2+1)^2}{144} \\
			=& k^2 {n \choose 2} (k+1)^2 \frac{k-1}{144}.
	\end{align*}	
	
	Sustituyendo en \ref{S-Friedman2} llegamos a 
	
	\[ \mathbb{E}[S] = \frac{nk(k^2-1)}{12}; \ 
		\var(S)= \frac{k^2n(n-1)(k+1)^2}{72} .\]
	
	Podemos definir la transformación
	
	\[ \chi^2_F = \frac{12S}{nk(k+1)} = 
	\frac{12 \sum\limits_{j=1}^k R_j^2}{nk(k+1)} -3n(k+1) .
	\]
	que tiene $\mathbb{E}[\chi^2_F] = k-1$, $\var(\chi^2_F)= 2(k-1)
\frac{n-1}{n} \approx 2(k-1)$, que son los dos primeros 
momentos de la distribución $\chi^2_{k-1}$. También los 
momentos de mayor orden de $\chi^2_F$ se aproximan a los 
momentos de mayor orden de $\chi^2_{k-1}$ para $n$ grande. A 
efectos prácticos, podemos tratar $\chi^2_F$ como una $
\chi^2_{k-1}$ para $n>7$ ó $n>10, k>5$. La región crítica $R$
para la hipótesis nula mencionada con nivel de significación 
$\alpha$ es
	\[ 
		\chi^2_F \in R \ \text{ para } 
			\chi^2_F \geq \chi_{k-1,\alpha}^2 .
	\]
	
\paragraph{Propuesta de Iman-Davenport} Sin embargo, este
test es demasiado conservador, por lo que Iman y Davenport
propusieron el estadístico
	\[
		F_F = \frac{(n-1)\chi^2_F}{n(k-1) - \chi^2_F},
	\]
	que sigue una distribución $F$ de Snedecor con $k-1$ y
$(k-1)(n-1)$ grados de libertad.
	
	
\subsubsection{Test de Friedman AR}
	Este test basado en el test de Friedman se propone como
solución a la hora de comparar un número pequeño de 
algoritmos para tener en cuenta el rendimiento no sólo 
para cada conjunto de datos sino entre ellos. El
siguiente algoritmo indica el procedimiento de cálculo.\\
	
\begin{algorithm}
	\caption{Cálculo del estadístico de Friedman-AR}
	\label{alg:Friedman-AR}
	\begin{algorithmic}[1]
	\REQUIRE $x_{ij}$, medida del rendimiento del
			algoritmo $j$ en el conjunto $i$.
	\STATE Se calcula la media del rendimiento de los 
		algoritmos para cada conjunto de datos, 
		$\bar{x}_{i \cdot}$
	\FOR {$ x_{ij}:\ i=1,\dots, n,\ j=1,\dots,k$}
		\STATE Observación alineada: $\leftarrow$
			Diferencia entre $x_{ij}$ y $\bar{x}_{i \cdot}$.	
	\ENDFOR
	\STATE Se ordenan de $1$ a $kn$ las observaciones 
		alineadas $R_{ij}$.
	\STATE Se calculan los totales de estos ránquines
		por algoritmo $\hat{R}_{\cdot j}$ y por
		conjunto de datos $\hat{R}_{i \cdot}$.
	\end{algorithmic}
\end{algorithm} 

	Entonces, se calcula el estadístico 

	\[
		T = \frac{(k-1)
				\left[
					\sum\limits_{j=1}^k
						\hat{R}_{\cdot j}^2 -
					kn^2(kn+1)^2/4
				\right]}
				{[kn(kn+1)(2kn+1)/6] 
				 - (1/k) \sum\limits_{i=1}^n 
				 			\hat{R}_{i \cdot}^2},
	\]
	que sigue una distribución $\chi^2_{k-1}$.
	
\subsubsection{Test de Quade}

	El test de Friedman considera todos los conjuntos de 
datos iguales en cuanto a importancia. Se propone la 
siguiente modificación modificando el peso en base a las
diferencias entre los rendimientos en cada base de datos.\\
	Para ello, se considera el rango de cada conjunto,
	\[ rango_i = \underset{j}{\max} \{x_{ij}\} - 
				 \underset{j}{\min} \{x_{ij}\},
	\]
	y se clasifican los rangos de menor a mayor, notando por 
$Q_i$ al puesto ocupado por el conjunto de datos $i$.
Entonces se multiplica $Q_i$ por la diferencia entre el 
ránquin dentro del conjunto $i$ ($R_{ij}$ y el puesto medio,
$\frac{k+1}{2}$:
	\[ 
		S_{ij} = Q_i 
				\left( R_{ij} - \frac{k+1}{2} \right).
	\]
	Denotamos por $S_j = \sum\limits_{i=1}^n S_{ij}$ y
calculamos los términos
\begin{align*}
		A_2 &= \frac{kn(n+1)(2n+1)(k+1)(k-1)}{72},\\
		B &= \frac{1}{n} \sum\limits_{j=1}^k S_j^2.
\end{align*}

	Finalmente, obtenemos el estadístico
	\[
		T_3 = \frac{(n-1)B}{A_2 - B}
	\]
	que sigue una distribución $F$ de Snedecor con $k-1$ y
$(k-1)(n-1)$.
	

\subsubsection{Procedimientos \textit{post-hoc}}

	Una vez se ha descartado la hipótesis nula de 
la equivalencia entre los clasificadores, se pretende
identificar dónde se producen las diferencias. Para ello,
no se considera una única hipótesis sino una familia 
de ellas consistentes en la igualdad del rendimiento entre 
algoritmos y pretendemos rechazarlas o no sin tener que ir 
una por una.\\
	Si consideramos la comparación de $k$ algoritmos con 
nivel de significación $\alpha$, la probabilidad de no 
cometer un error de tipo I es $1-\alpha$, y entonces la 
probabilidad de no cometerlo en las $k-1$ comparaciones es
$(1-\alpha)^{k-1}$. Así pues, la probabilidad de cometer un
error de tipo I es $1-(1-\alpha)^{k-1}$, lo que aumenta 
la probabilidad de cometer un error cuando crece el número 
de algoritmos a comparar. Para ello, la solución es 
ajustar los $p$-valores (\textit{Adjusted $p$-values}, APV), 
y realizar las comparaciones con $\alpha$. Se describen a continuación algunos procedimientos para ajustar los 
$p$-valores y los valores de $\alpha$:

\begin{itemize}
\item Procedimientos de un paso.
	\begin{description} 
	\item[Bonferroni-Dunn] Ajusta el valor de $\alpha$ 
		dividiéndolo por el número de comparaciones
		realizadas, $k-1$. $APV_i = \min \{\nu,1\}$,
		con $\nu = (k-1)p_i$.
	\end{description}
\item Procedimiento descendente.
	\begin{description} 
	\item[Holm] Consideramos los $p$-valores ordenados
		de manera ascendente $p_1 \leq \dots \leq p_{k-1}$
		y $H_1, \dots, H_{k-1}$ sus hipótesis asociadas.
		Se rechazan las hipótesis $H_1, \dots, H_{i-1}$
		para $i$ el menor entero tal que $p_i > 
		\alpha/(k-i)$. $APV_i = \min \{\nu,1\}$,
		con $\nu = \max \{(k-j)p_j: 1 \leq j \leq i\}$.
	\item[Holland] Partiendo de la misma situación que en
		el de Holm, se rechazan las hipótesis $H_1, \dots, 
		H_{i-1}$ para $i$ el menor entero tal que $p_i > 
		1- (1-\alpha)^(k-i)$. $APV_i = \min \{\nu,1\}$,
		con $\nu = \max \{1- (1-\alpha)^(k-j): 1 \leq j 
		\leq i\}$.
	\item[Finner] Partiendo de la misma situación que en
		los anteriores, se rechazan las hipótesis 
		$H_1, \dots, H_{i-1}$ para $i$ el menor entero tal 
		que $p_i >  1- (1-\alpha)^{(k-1)/i}$. 
		$APV_i = \min \{\nu,1\}$, con $\nu = \max \{1- 
		(1-\alpha)^{(k-1)/j}: 1 \leq j \leq i\}$.
	\end{description}
\item Procedimiento ascendente.
	\begin{description} 
	\item[Hochberg] Se ajusta el valor de $\alpha$ de manera
		ascendente. Se compara el mayor $p$-valor con 
		$\alpha$, el siguiente con $\alpha/2$, el siguiente
		con $\alpha/3, \dots$ hasta encontrar una hipótesis
		que se pueda rechazar, rechazando también aquellas
		con un $p$-valor menor. $APV_i = \max \{(k-j)p_j:
		(k-1) \geq j \geq i \}$.
	\item[Hommel] Este método es más complejo de calcular
		y comprender. Se busca el mayor $j$ tal que 
		$p_{n-j+k} < k\alpha/j$ para todo $k=1,\dots,j$.
		Si no existe $j$, se rechazan todas las hipótesis,
		si no, se rechazan todas para las que $p_i \leq 
		\alpha/j$.
	\item[Rom] Consiste en una modificación del de Hochberg,
		calculando $\alpha$ de la siguiente manera:
		\[ 
			\alpha_{k-i} = \frac{
								\sum\limits_{j=1}^{i-1}
									\alpha^j
								- \sum\limits_{j=1}^{i-2}
									{i \choose k}
									\alpha_{k-1-j}^{i-j}
						}{i},
		\]
		con $\alpha_{k-1} = \alpha, \alpha_{k-2}= \alpha/2$. 
		$APV_i = \max \{(r_{k-j})p_j: (k-1) \geq j \geq 
		i \}$, con $r_{k-j}=\alpha/\alpha_{k-j}$
	\end{description}	
\item Procedimiento en dos pasos.
	\begin{description}
	\item[Li] En primer lugar se rechazan todas las 
		hipótesis si $p_{k-1} \leq \alpha$. Si no, 
		no se rechazará $H_{k-1}$. En segundo lugar, se 
		rechazan las hipótesis $H_{i}$ tal que 
		$p_i \leq \frac{1-p_{k-1}}{(1-\alpha)\alpha}$.
		$APV_i = p_i/(p_i + 1 - p_{k-1})$.	
	\end{description}	 
\end{itemize}

\section{Test basados en permutaciones}
	
	En este apartado se incluye una descripción general sobre 
los test basados en permutaciones. La hipótesis nula a la 
hora de realizar inferencia es considerar que una muestra 
aleatoria $\mathbf{x}$ proviene de una misma población 
$P \in \mathcal{P}$. Por tanto, supuesta esta hipótesis nula, 
obtenemos que si la muestra $\mathbf{x}$ viniese de 
poblaciones distintas, los valores pertenecientes a cada 
población podrían intercambiarse por los de otra sin ningún 
problema.\\
	Notemos además que $\mathbf{x}$ es un conjunto de 
estadísticos suficientes para cualquier distribución en 
$H_0$. Esto se debe a que, si $H_0$ es cierto, sea $f_P$ la 
densidad de $P$, $f_P^{(n)}(x)$ es la densidad de la variable 
$X^{(n)}$. Como $f_P^{(n)}(\mathbf{x}) = f_P^{(n)}(\mathbf{x}) 
\cdot 1$ para todo $\mathbf{x} \in \Omega^n$, excepto para 
$f_P^{(n)}(\mathbf{x}) = 0$ debido al teorema de 
factorización (\cite[p. 167]{ibarrola1993principios}) cualquier conjunto $\mathbf{x}$ 
es un conjunto 
de estadísticos suficientes para cualquier $P \in 
\mathcal{P}$.\\
	En estos test se realizan inferencias condicionadas con
respecto a la muestra $\mathbf{x}$, que es a su vez un conjunto de 
estadísticos suficientes en $H_0$. Este hecho unido a la 
intercambiabilidad de los datos observados con respecto a los 
grupos hace que los test sean independientes de las 
distribución $P$, por lo que no necesitamos conocerla. Esto 
queda más claramente explicado en el siguiente principio:

\begin{definicion}[Principio de los test basados en 
permutaciones]
	Si dos experimentos toman valores en el mismo espacio
muestral $\Omega$ respectivamente con las distribuciones 
$P_1, P_2 \in \mathcal{P}$, dan el mismo conjunto de datos, 
entonces las inferencias condicionales a los datos obtenidas 
usando el mismo test estadístico debe 
ser la misma, debido a que la intercambiabilidad de los datos con 
respecto a los grupos se cumple en la hipótesis nula. Por 
consiguiente, si dos experimentos con distribuciones 
$P_1, P_2$ dan respectivamente como muestras $\mathbf{x_1}, 
\mathbf{x_2}$, con $\mathbf{x_1} \neq \mathbf{x_2}$, entonces 
las inferencias con respecto a $\mathbf{x_1}$ y 
$\mathbf{x_2}$ pueden ser distintas.
\end{definicion}


\subsection{Ejemplo con muestras apareadas}

	Para explicar el funcionamiento de los test basados en 
permutaciones, expondremos situaciones ya conocidas en las 
que sabemos cómo utilizar test paramétricos o test no 
paramétricos expuestos anteriormente. El primer caso es el de 
muestras pareadas. Supondremos una situación en la que 
tenemos 20 problemas y comparamos el rendimiento de dos
algoritmos en cada uno de ellos.
Por tanto consideramos la v.a. $Y = (Y_1, Y_2)$ de $n=20$ 
elementos $\{(Y_{1i},Y_{2i}): i = 1, \dots, n\}$. La 
hipótesis nula es que las diferencias en que las 
distribuciones marginales de $Y_1$ e $Y_2$ son iguales:
$H_0 : \{ Y_1 = Y_2 \}$ frente a $H_1: \{ Y_1 > Y_2 \}$\\
	
	En primer lugar, nótese que la hipótesis nula $H_0$ 
implica que las dos variables $Y_1, Y_2$ son intercambiables. 
Esto es, supuesto $H_0$, para cada individuo los dos valores 
observados podrían corresponder con cualquiera de las 
observaciones, para un algoritmo u otro. Es decir, si 
consideramos la variable aleatoria $X = Y_2 - Y_1$, cada 
valor observado $X_i, i=1, \dots, n$ podría tener cualquier 
signo con probabilidad $1/2$. Una forma de realizar el test 
sería considerando el estadístico $T = \sum\limits_{i=1}^n 
X_i$, cuya distribución condicionada a $\mathbf{X} = \{X_i: i 
= 1, \dots, n\}$, $F_T(t|\mathbf{X})$ vendrá dada bajo $H_0$ 
considerando la atribución de todas las formas posibles del 
signo de cada diferencia con probabilidad $1/2$. Denotamos
por $T^* = \sum\limits_{i=1}^n X_i^*$, donde 
$X_i^*$ se obtiene asignando aleatoriamente el signo. 
Notaremos por $\Omega_{/\mathbf{X}}$ al espacio muestral de 
permutaciones. Este espacio consiste en las posibles 
permutaciones de la muestra obtenidas. En este caso, el espacio 
tiene $M^{(n)} = 2^v$ elementos, donde $n-v$ es el número de 
diferencias nulas. Notamos por $F(t|\mathbf{X}) = P[T^* \leq 
t | \mathbf{X}], t \in \mathbb{R}$ a la cdf condicional de la 
permutación inducida por $T$ dado $\mathbf{X}$.\\
	Si el tamaño de la muestra $n$ es grande, podrá aplicarse 
el Teorema Central del Límite Permutacional (\ref{th:PCTL}) 
para aproximar la distribución de la permutación 
$F[t|\mathbf{X}]$. Para resolver un problema de este tipo 
podemos guiarnos por lo siguiente:
	
\begin{enumerate}[a]
	\item Si $n \leq 25$ es posible calcular exactamente 
		$T^*$  para todas las posibles permutaciones y 
		calcular $F[t|\mathbf{X}]$.
	\item Si $n \geq 200$, $\sigma_X$ se supone finito y el 
		cociente $(\sum_i X_i^4)/(\sum_i X_i^2)^2$ es 
		pequeño, $F[t|\mathbf{X}]$ puede aproximarse por 
		el PCTL (\ref{th:PCTL}). 
	\item En otro caso, $F[t|\mathbf{X}]$ puede aproximarse 
		utilizando el algoritmo condicional de Monte 
		Carlo (CMC) (\ref{alg:CMC-pvalue}).
\end{enumerate}
	
\subsubsection{Teoría de los test basados en permutaciones}

\begin{definicion}[Espacio de referencia condicional]
	Conjunto de puntos del espacio muestral $\Omega^n$ que 
son equivalentes a $\mathbf{X}$ en términos de la información 
asociada a la verosimilitud subyacente. Lo notamos por 
$\Omega^n_{/\mathbf{X}}$.
\end{definicion}

	Esto es, $\Omega^n_{/\mathbf{X}}$ contiene todos los 
puntos $\mathbf{X}^*$ tales que $dP(X)/dP(X^*)$ es 
independiente de $P$. Dado que en $H_0$ la probabilidad de 
obtener $\mathbf{X}$ es la misma que de obtener una 
permutación $\mathbf{X}^*$,
	\[ 
		\Omega^n_{/\mathbf{X}} = 
			\left\lbrace
				\underset{\mathbf{u}^*}{\bigcup}
					[ X(u_i^*), i = 1, \dots, n ]
			\right\rbrace ,
	\]
con $\mathbf{u}^*$ cualquier permutación de las etiquetas 
$(1, \dots, n)$.

\paragraph{Definicion de los test basados en permutaciones}
\begin{definicion}[Soporte de la permutación]
	Para realizar un test estadístico es necesaria una 
función no degenerada $T: \Omega^n \rightarrow \mathbb{R}$. 
Definimos el soporte de la permutación de $(T, \mathbf{X})$ 
como $\mathcal{T}_\mathbf{X} = \{ T^* = T(\mathbf{X}^*): 
\mathbf{X}^* \in \Omega_{/\mathbf{X}}\}$ 
\end{definicion}	

	Supongamos $H_0$ cierta. Entonces, $\mathbf{X}^*$ se 
distribuye uniformemente sobre $\Omega_{/\mathbf{X}}$. 
Pongamos los $M^{(n)}$ elementos de $\mathcal{T}_\mathbf{X}$ 
en orden creciente. De esta manera para $\alpha \in (0,1)$, 
se define $T_\alpha(\mathbf{X}) = T_\alpha = T^*_{(M^{(n)}_
\alpha)}$ como el valor crítico asociado a $(T, \mathbf{X})$, 
donde $M^{(n)}_\alpha$ es el número de valores $T^*$ 
estrictamente menor que $T_\alpha$. Nótese que $T_\alpha$ 
depende de $\Omega_{/\mathbf{X}}$ y no únicamente de $
\mathbf{X}$, pues $T_\alpha(\mathbf{X}) = T_
\alpha(\mathbf{X}')$ para $\mathbf{X}' \in \Omega_{/
\mathbf{X}}$. Esto implica que si $\alpha$ es fijo,$T_\alpha$ 
es fijo para $\mathcal{T}_\mathbf{X}$.
	
\paragraph{Test aleatorizados basados en permutaciones }

	Se define la versión aleatorizada del test basado en 
permutaciones para $(T, \mathbf{X})$ como 
	\[ 
		\phi_R = \left\lbrace \begin{array}{cc}
			1 		& \text{ si } T^0 > T_\alpha \\
			\gamma 	& \text{ si } T^0 = T_\alpha \\
			0		& \text{ si } T^0 < T_\alpha \\
		\end{array} \right. ,
	\]
	donde $T^0 = T(\mathbf{X})$, el valor de $T$ en los datos 
observados y $\gamma$ es:
	\[
		\gamma = \frac{\alpha - P[T^0 > T_\alpha | 
								\Omega_{/\mathbf{X}}]}
					  {P[T^0 = T_\alpha | 
								\Omega_{/\mathbf{X}}]} .
	\]
	
	Para aplicar $\phi_R$ es habitual usar un experimento 
aleatorio independiente de $\mathbf{X}$. Por ejemplo, 
rechazando $H_0$ si $U \geq \gamma$ para $T^0 = T_\alpha$, 
con $U$ un valor aleatorio de la variable uniforme 
$\mathcal{U}(0,1)$.\\
	Para $\mathbf{X} \in \Omega^n$ la esperanza en $H_0$ de 
$\phi_R$ para $\alpha$ es:
	\[ 
		\mathbb{E}[ \phi_R(\mathbf{X}) | \Omega_{/\mathbf{X}}] =
			P[ T^0(\mathbf{X}) > T_\alpha(\mathbf{X}) | 
								\Omega_{/\mathbf{X}}] + 
			P[ T^0(\mathbf{X}) = T_\alpha(\mathbf{X}) | 
								\Omega_{/\mathbf{X}}] =
			\alpha .
	\]
	
	Por consiguiente, el tamaño del test aleatorizado basado en 
permutaciones es exactamente $\alpha$.
	
\begin{proposicion}
	Supuesta la condición de intercambiabilidad para 
$\mathcal{X}$. Entonces para todas las distribuciones $P$ y 
uniformemente para todos los conjuntos de datos $\mathbf{X} 
\in \Omega^n$ la probabilidad condicionada de rechazo de 
$\phi_R$ es invariante con respecto a $\mathbf{X}$ y $P$.
\end{proposicion}

\paragraph{Test basados en permutaciones no aleatorizados}

	En los contextos de aplicación se suele utilizar la 
versión no aleatorizada. Se define
	
	\[ 
		\phi = \left\lbrace \begin{array}{cc}
			1 		& \text{ si } T^0 \geq T_\alpha \\
			0		& \text{ si } T^0 < T_\alpha \\
		\end{array} \right. .
	\]
	
	Ahora, el error asociado de tipo I es 
	
	\[ 
		\mathbb{E}[ \phi(\mathbf{X}) | \Omega_{/\mathbf{X}}] =
		P[ T^0 \geq T_\alpha | \Omega_{/\mathbf{X}}] =
		\sum\limits_{\Omega_{/\mathbf{X}}}
			\mathbb{I}[ T(\mathbf{X}^*) \geq T_\alpha ]/
				M^{(n)} =
		\alpha_a \geq
		\alpha,
	\] 
	
	donde $\alpha_a$ es el llamado $\alpha$-valor alcanzable 
asociado a $(T, \mathbf{X})$. Para $(T,\mathbf{X})$, el 
$\alpha$-valor alcanzable pertenece a $\Lambda_\mathbf{X}
^{(n)} = \{ L_\mathbf{X}(t): dL_\mathbf{X}(t) > 0 \}$, para 
la función $ L_\mathbf{X}(t) = P[ T^* \geq t | \Omega_{/
\mathbf{X}}] $. 
	
\paragraph{$p$-Valor} Determinar el valor crítico $T_\alpha$ 
para un test estadístico $T$ puede ser complicado en la 
práctica. Por tanto, normalmente se hace referencia a un $p$-
valor asociado a $(T, \mathbf{X})$, que se define como $p = 
\lambda_T(\mathbf{X}) = L_\mathbf{X}(T^0) = P[ T^* \geq T^0 | 
\Omega_{/\mathbf{X}}]$. Este valor se puede calcular 
enumerando completamente $\mathcal{T}_\mathbf{X}$ o 
estimándolo mediante el algoritmo condicional de Monte Carlo. 
El $p$-valor $p$ es una función creciente de $T^0$ y tiene 
una  correspondencia uno a uno con los $\alpha$-valores 
alcanzables de $\phi$, puesto que $\lambda_T(\mathbf{X}) > 
\alpha$ implica $T^0 < T_\alpha$ y viceversa.

\paragraph{Algoritmo CMC para estimar el $p$-valor}  
Describimos a continuación los pasos generales de un 
algoritmo condicional de Monte Carlo para la evaluación de un 
$p$-valor de un test estadístico $T$ en un conjunto de datos 
$\mathbf{X} = \{ X_i, \ i=1, \dots, n; n_1, n_2 \}$

\begin{algorithm}
	\caption{Algoritmo CMC para estimar el $p$-valor}
	\label{alg:CMC-pvalue}
	\begin{algorithmic}[1]
	\REQUIRE
		\begin{enumerate}[a]
		\item Muestra $\mathbf{X}$
		\item Estadístico $T$
		\item Número de iteraciones $B$
		\end{enumerate}
		\STATE Calcular $T^0 = T(\mathbf{X})$
		\FOR { $b \in \{1, \dots, B\}$ }
			\STATE Realizar una permutación $\mathbf{X}^*$ de  
			$\mathbf{X}$.
			\STATE Calcular $T^* = T(\mathbf{X}^*)$
		\ENDFOR
		\STATE El conjunto $\{ \mathcal{X}^*_b, b = 1, \dots, 
		B \}$ es una muestra aleatoria de 
		$\Omega_{/\mathbf{X}}$ y los valores correspondientes 
		$\{ T^*_b, b = 1, \dots, B \}$ simulan la 
		distribución nula de las permutaciones de $T$ 
		(aquella que tendría $T$ de ser cierta $H_0$). 
		Entonces el $p$-valor se estima como $\hat{\lambda}
		(\mathbf{X}) = \sum\limits_{b=1}^B 
		\mathbb{I}[T_b^* \geq T^0]/B$, esto es, la proporción 
		de los valores de la permutación mayores que el 
		observado.
	\end{algorithmic}
\end{algorithm}

\paragraph{Teorema Central del Límite Permutacional}
	Por último en esta sección, se presenta la versión 
permutacional del TCL de Wald y Wolfowitz (\cite{wald1944}).
Notaremos por $\mathbf{X} = \{X_1, \dots, X_n\}$ de $n$ v.a. 
independientes e idénticamente distribuidas. Supongamos que 
la hipótesis nula sobre la que se realiza el test, $H_0$, 
implica que las permutaciones en $\mathbf{X}$ son igualmente 
probables.\\
	Definimos el test estadístico permutacional
	\[
		T = T(\mathbf{X}) = 
			\sum\limits_{i=1}^n
				A_i X_i,
	\]
	para el vector $\mathbf{A} = \{A_1, \dots, A_n\}$. Para 
un vector genérico $\mathbf{D}$, notamos $\bar{D} = 
\frac{\sum\limits_{i=1}^n D_i}{n}$, $\mu_r(\mathbf{D}) = 
\frac{\sum\limits_{i=1}^n (D_i - \bar{D})^r}{n}$, 
$W_r(\mathbf{D}) = \frac{\sum\limits_{i=1}^n 
|D_i - \bar{D}|^r}{n}$, $R(\mathbf{D}) = 
\underset{1 \leq i \leq n}{\max} D_i - 
\underset{1 \leq i \leq n}{\min} D_i$.\\

	De la distribución de la permutación de $T$ se obtiene

\begin{align*}
	\mathbb{E}_{\Omega_{/\mathbf{X}}} \left[
		T(\mathbf{X}^*) | \Omega_{/\mathbf{X}}
		\right] &=
		n \bar{X} \bar{A}, \\
	\mathbb{E}_{\Omega_{/\mathbf{X}}} \left[
		(T(\mathbf{X}^*) - n\bar{X}\bar{A})^2 | 
			\Omega_{/\mathbf{X}}
		\right] &=
		\frac{[n^2 \mu_2(\mathbf{X}) \mu_2(\mathbf{A})]^2}
			{n-1}.
\end{align*}

	Entonces, escribimos la v.a.
	\[
		Z = \frac{(T - n \bar{X} \bar{A})\sqrt{n-1}}
				{n^2 \mu_2(\mathbf{X}) \mu_2(\mathbf{A})},
	\]
	y sabemos que la cdf de $Z$, $F_Z(t|\Omega^n)$ puede 
determinarse enumerando todas las permutaciones de 
$\mathbf{X}$. Antes de enunciar el teorema, enunciaremos la 
condición sobre $\mathbf{X}$ y $\mathbf{A}$. 
$\{ \mathbf{D}_n; n \geq 2 \}$:

\begin{equation}
	 \label{eq:PCLTcond}
	 \tag{C.1}
	 \frac{\mu_r(\mathbf{D}_n)}
	 		{\mu_2(\mathbf{D})^{r/2}} = O(1),\quad
	 		\forall r \in \mathbb{N},\ r \geq 2
\end{equation}

\begin{teorema}[Teorema Central del Límite Permutacional]
	Sean $\{\mathbf{A}_n; n \geq 2\}$, $\{\mathbf{X}_n; 
n \geq 2\}$ cumpliendo la condición~\ref{eq:PCLTcond}. 
Entonces para $n \rightarrow \infty$,
	\[ F_Z(t | \Omega^n_{\mathbf{X}}) 
		\overset{P}{\rightarrow} \mathit{N}(t | 0,1)
	\]
	\label{th:PCTL}
\end{teorema}


\chapter{Test bayesianos}	
\label{chapter:test_bayesianos}

	En esta sección se hace una introducción a conceptos de 
la estadística bayesiana y test basados en esta estadística
para su aplicación en la comparación de algoritmos de
aprendizaje automático.\\
	Según Gelman (\cite{Gelman03-105949}), la inferencia 
bayesiana es el proceso de ajustar un modelo de probabilidad 
a un conjunto de datos y resumir el resultado  mediante una 
distribución de probabilidad sobre los  parámetros del modelo 
y sobre cantidades no observadas como predicciones para 
nuevas observaciones. Estamos por tanto frente a una 
forma de entender la estadística distinta, cuya principal
característica es el uso de la probabilidad para 
cuantificar la incertidumbre en las inferencias basadas en
el análisis de los datos. Este análisis puede dividirse en
tres pasos:

	\begin{enumerate}
	\item establecimiento de un modelo de probabilidad;
	\item condicionamiento sobre los datos observados,
		calculando e interpretando la distribución 
		\textit{a posteriori};
	\item evaluación del modelo y las implicaciones de la
		distribución \textit{a posteriori} resultante.
	\end{enumerate}

\section{Problemas de los test de hipótesis nula}	

	Los test de hipótesis nula (THN) son en ocasiones como un
método infalible para la realización de un experimento y la
forma de obtener conclusiones. Sin embargo, existen críticas
hacia el uso de este tipo de test, principalmente por parte 
de estadísticos bayesianos, causadas en su mayoría por el mal
uso que se hace de ellos. Las críticas realizadas son las
siguientes:

\begin{itemize}
\item Las decisiones dicotómicas. Al fijar el valor 
	$\alpha$ (normalmente $0.05$) se pueden realizar 
	decisiones opuestas con una mínima variación del $p$-
	valor. Se suele también confundir el significado del 
	test al considerar equivalentes dos métodos para los que
	no se encuentren diferencias estadísticamente 
	significativas.
\item Los THN no estiman la probabilidad
	de la hipótesis. La pregunta que habitualmente nos
	hacemos al comparar por ejemplo dos clasificadores es
	``\textit{¿cuál es la probabilidad de que el rendimiento
	de dos clasificadores sea el mismo?}'', en términos
	formales, $P(H_0 | \mathbf{X})$. La pregunta que
	responden es $P(T(\mathbf{X}) > \tau_\alpha | H_0)$.
\item Los test de hipótesis nula sobre puntos son
	prácticamente siempre falsos. La diferencia entre
	dos clasificadores puede ser muy pequeña, sin embargo
	no hay dos clasificadores cuyo rendimiento sea 
	equivalente. Esto tiene como consecuencia que añadiendo
	suficientes puntos es posible encontrar diferencias
	significantes. Además, no sólo en este contexto, los
	experimentos realizados y publicados suelen considerar
	hipótesis nulas difícilmente sostenibles donde el
	investigador busca rechazarla, con lo que el test está
	resulta sesgado.
\item El $p$-valor no distingue entre el efecto del tamaño
	y el tamaño de la muestra. Si el efecto del tamaño de
	$H_0$ es pequeño, se necesita más datos para demostrar 
	la diferencia. Sin embargo, como el tamaño de la muestra
	depende del investigador, y por el punto anterior $H_0$ 
	``es siempre falsa'', con suficientes datos se puede
	rechazar la hipótesis nula. En cambio, si se dispone 
	de pocos datos, se puede fallar al no rechazar una
	hipótesis nula ampliamente falsa.
\item Los THN no tienen en cuenta la magnitud y la 
	incertidumbre. No se proporciona información sobre la 
	magnitud del efecto o la incertidumbre de la estimación
	realizada, con lo que se dificulta un posterior análisis
	de las propias conclusiones realizadas.
\item No se obtiene información cuando no se rechaza la
	hipótesis nula. Es un error habitual al considerar 
	que al no encontrar diferencias significativas se
	considere cierta la hipótesis nula. 
\item No hay manera de decidir previamente el nivel 
	de significación $\alpha$. Aunque habitualmente se 
	fija un nivel de significación de $\alpha = 0.05$, y es
	este valor el que determina los límites para rechazar 
	$H_0$, este valor no tiene un significado concreto para
	el experimento. 
\item La inferencia depende del propósito del estudio. 
	La distribución de la muestra cambia por ejemplo si
	se comparan dos algoritmos en todos los conjuntos de
	datos disponibles que si se trata de realizar un 
	número de observaciones en un único conjunto.
\end{itemize}
	
\section{Análisis de los resultados experimentales}

	Hay dos enfoques para el análisis de resultados 
experimentales. El primero, conocido como \textbf{factor de 
Bayes} o \textbf{comparación de modelo bayesiano}, está, como 
los THN, basado en lo que llamaríamos un valor nulo. Se 
comparan dos modelos y se escoge el más probable según los 
datos obtenidos. Éste modelo también recibe algunas críticas
similares a los THN pues en cierta forma se está cambiando el
valor $\alpha$ por el usado para elegir entre un modelo u 
otro, volviendo a una decisión dicotómica. El otro modelo, el 
cual se desarrollará con más detalle, se conoce como 
\textbf{estimación bayesiana de parámetros}. Al realizar
el análisis sólo hay que establecer un rango de valores 
candidatos para el parámetro (modelo \textit{a priori}) 
y usar la inferencia bayesiana para calcular la probabilidad
de todos los valores candidatos (distribución \textit{a 
posteriori}).\\

\section{t-test bayesiano correlado}
	\label{ssec:bayes-ttest}
	Este test se usa para el análisis de la validación
cruzada en un único conjunto de datos. Su aporte consiste en 
tener en cuenta la correlación debida al solapamiento de los 
conjuntos de entrenamiento. El modelo de los datos 
considerados es
	\[ \mathbf{X} = \mu \mathbf{1} + \mathbf{v}, \]
	con $\mathbf{X} = (X_1, \dots, X_n)$, $\mathbf{1}$ un 
vector de unos de $n$ elementos, $\mu$ el parámetro de
interés, esto es la media de las distancias entre el 
rendimiento de los clasificadores y $\mathbf{v} \backsim 
\Sigma_{n \times n}$, una normal multivariante con
media cero y matriz de covarianzas $\Sigma$. Se caracteriza
$\Sigma$ de la siguiente manera: $\Sigma_{ii} = \sigma^2$
y $\Sigma_{ij} = \sigma^2 \rho$ para $i \neq j$, con $\rho$
la correlación y $\sigma^2$ la varianza. Entonces,
la probabilidad del modelo es 
	\[ 
		P(\mathbf{X} | \mu, \Sigma) = 
		\frac{\exp\left(
				-\frac{1}{2}
				(\mathbf{X}-\mu\mathbf{X})^T
				\Sigma^{-1}
				(\mathbf{X}-\mu\mathbf{X})
			\right)}
		{(2\pi)^{n/2} \sqrt{|\Sigma|}}.
	\]

	Esta probabilidad no permite estimar $\rho$ de los datos, 
así que adoptamos la heurística $\rho = \frac{n_{test}}{n}$,
con $n_{test}$ el número de datos en el conjunto de test.
Pretendemos estimar los parámetros $\mu$ y $\nu = 
\frac{1}{\sigma^2}$, para lo cual consideramos la 
distribución \textit{a priori} Normal-Gamma
	\[
		P(\mu, \nu | \mu_0, k_0, a, b) =
			N\left(\mu; \mu_0, \frac{k_0}{\nu} \right)
			G(\nu; a, b) = 
			NG(\mu, \nu ; \mu_0, k_0, a, b).
	\]
	Escogiendo los parámetros $\mu_0 = 0, k_0 \rightarrow 
\infty, a = -\frac{1}{2}, b=0$ la distribución \textit{a 
posteriori} de $\mu$ resulta ser la distribución de Student:
	
	\[
		P(\mu | \mathbf{X}, \mu_0, k_0, a, b) =
			St \left(
					\mu; n-1 , \bar{X},
					\left( \frac{1}{n} + 
							\frac{\rho}{1 - \rho} 
					\right) \hat{\sigma}^2
			   \right),
	\]
	con $\bar{X}$ y $\hat{\sigma}^2$ los estimadores 
habituales. Por tanto podemos hacer el análisis sobre $\mu$ 
basándonos en esta distribución.


\paragraph{Ejemplo} Supongamos que realizamos un experimento
con este test, realizando la comparación entre dos algoritmos
para una base de datos. La información obtenida es una 
distribución $T$ de Student, que describe como se ha mencionado
la ditribución de la media de las diferencias del rendimiento
entre dos clasificadores. Ahora, evaluamos la probabilidad 
de la hipótesis que queramos realizar, por ejemplo, que un 
algoritmo es mejor que otro directamente en la distribución
de $\mu$, calculando la probabilidad de que $\mu > 0$ ó $\mu 
< 0$ según corresponda. Para decir que los dos algoritmos son
prácticamente equivalentes, podemos suponer que esto 
significa que la diferencia de la precisión es menor que un
$1\%$. El intervalo $[-0.01, 0.01]$ se conoce entonces como
\textbf{región de práctica equivalencia} (\textit{rope}). Una manera 
de mostrar la incertidumbre de la estimación es indicando el 
intervalo de los valores que son más confiables y cubren el 
$q\%$ de la distribución (HDI). 

\paragraph{Toma de decisiones} A la hora de realizar 
análisis queremos obtener conclusiones, para lo que nos 
apoyaremos en la \textit{rope}. Si tenemos dos algoritmos $f_1, f_2$, 
notaremos por $P(f_1 \ll f_2)$ a la probabilidad de que el 
primero tenga un rendimiento inferior, por $P(f_1 \gg f_2)$ a 
que la de que tenga un rendimiento superior y por $P(f_1 =
f_2)$ a la probabilidad de que sean prácticamente 
equivalentes. Podemos tomar las decisiones con un nivel
de confianza (ahora con una interpretación directa) si
tenemos una de estas tres situaciones con un valor mayor
al nivel de confianza deseado.

\section{Test bayesiano de signo y \textit{signed-rank}}

	Al igual que el $t$-test, el test no paramétrico de 
signo también cuenta con la versión bayesiana. Partimos
del vector $z = \{z_1, \dots, z_q\}$ de muestras i.i.d.. 
Suponemos como distribución \textit{a priori} un proceso
de Dirichlet ($DP$). Un $DP$ es una distribución sobre 
distribuciones de probabilidad tales que sus distribuciones
marginales sobre particiones finitas siguen la distribución
de Dirichlet (que se puede considerar la generalización
multivariante de la distribución Gamma). El $DP$ está
definido por dos parámetros: la fuerza \textit{a priori},
$s>0$ y la media \textit{a priori} , $G_0$. Tengamos en 
cuenta que al ser $DP$ una distribución sobre distribuciones,
$G_0$ es una medida de probabilidad. Escogiendo $G_0 = 
\delta_{z_0}$, una delta de Dirac centrada en el punto $z_0$,
la densidad de probabilidad \textit{a posteriori} de $Z$ tiene la
expresión:
	\[ P(z) = 
		w_0 \delta_{z_0}(z) +
			\sum\limits_{j=1}^n
				w_j \delta_{z_j}(z); \quad
		(w_0, w_1, \dots, w_n) \sim Dir(s,1, \dots, 1).
	\] 
	
	Si observamos esta expresión, tenemos que es una 
combinación de deltas de Dirac centradas en las observaciones
y en la observación \textit{a priori} $z_0$ y cuyos pesos
siguen una distribución de Dirichlet. Para el análisis
experimental procedemos de la siguiente manera. Calculamos:

\begin{align*}
	\theta_l &= P[ z < -r ] = 
				\sum\limits_{i = 0}^q 
					w_i \mathbb{I}[ z_i \in (-\infty,-r)],\\
	\theta_e &= P[ |z| \leq r ] = 
				\sum\limits_{i = 0}^q 
					w_i \mathbb{I}[ z_i \in [-r,r]], \\
	\theta_r &= P[ z > r ] = 
				\sum\limits_{i = 0}^q 
					w_i \mathbb{I}[ z_i \in (r, \infty)].
\end{align*}

	Así, obtenemos una versión bayesiana del test de signo, 
que tiene en cuenta la \textit{rope} $[-r,r]$. De hecho, $\theta_l,
\theta_e, \theta_r$ son las probabilidades de que la media
de las diferencias esté en el respectivo intervalo. 
Debido a que $(w_0, w_1, \dots, w_n) \sim Dir(s,1, \dots, 
1)$, se demuestra que
	\[ 
		\theta_l,\theta_e, \theta_r \sim
		Dir(n_l + s \mathbb{I}[ z_0 \in (-\infty,-r) ],
			n_e + s \mathbb{I}[ z_0 \in [-r,r] ],
			n_r + s \mathbb{I}[ z_0 \in (r, \infty) ])
	\]
	donde $n_l, n_e, n_r$ son el número de observaciones 
que caen en cada intervalo. En definitiva, para 
completar el modelo, debemos escoger el valor \textit{a 
priori} de la fuerza $s$ y la observación $z_0$.\\

	Si ahora consideramos
	
\begin{align*}
	\theta_l &= P[ z < -r ] = 
				\sum\limits_{i = 0}^q 
				\sum\limits_{j = 0}^q 
					w_i w_j 
					\mathbb{I}[ z_i + z_j 
								\in (-\infty,-2r)],\\
	\theta_e &= P[ |z| \leq r ] = 
				\sum\limits_{i = 0}^q 
				\sum\limits_{j = 0}^q 
					w_i w_j 
					\mathbb{I}[ z_i + z_j 
								\in [-2r,2r]], \\
	\theta_r &= P[ z > r ] = 
				\sum\limits_{i = 0}^q 
				\sum\limits_{j = 0}^q 
					w_i w_j 
					\mathbb{I}[ z_i + z_j 
								\in (2r, \infty)],
\end{align*}
	obtenemos una versión bayesiana del 
test de rangos con signo, con la \textit{rope} $[-r,r]$. 
Ahora la distribución de 
$\theta_l, \theta_e, \theta_r$ no está claramente definida,
pero se puede calcular muestreando mediante el método
de Monte Carlo los pesos $(w_0, w_1, \dots, w_n) \sim 
Dir(s,1, \dots, 1)$.\\
	Para ambos métodos falta por fijar $s$ y $z_0$. 
En los experimentos realizados por Benavoli 
(\cite{DBLP:journals/corr/BenavoliCDZ16}) se utiliza $s=0.5$ y $z_0=0$. Se realizan 
también experimentos con $z_0 = \infty$ y $z_0 = -\infty$, 
obteniendo resultados muy parecidos en la probabilidad
de escoger un algoritmo frente a otro.




%
\part{Desarrollo informático}
\label{part:informatica}

\chapter{Revisión del estado del arte}
\label{chapter:revision}

\section{Test paramétricos}

\paragraph{Problemas a evitar y enfoque recomendado} 
\cite{Salzberg1997} Este es uno de los primeros artículos 
que hace una serie de recomendaciones para la realización
de experimentos en el ámbito de la comparación 
de algoritmos de aprendizaje automático. Se incluyen 
también varias consideraciones sobre los problemas y
prácticas incorrectas al realizar estos experimentos. Los
test propuestos para la comparación son el test de McNemar, 
el $t$-test y el test ANOVA para múltiples algoritmos. Para
la realización de varios experimentos independientes 
se propone el ajuste de Bonferroni del nivel de significación
$\alpha$, que necesita que los experimentos sean
independientes.


\subsection*{5$\times$2 CV $t$-test}
	
	En la introducción (\ref{sec:CV}), se describe el 
proceso de validación cruzada, especialmente útil cuando se
disponen de pocos datos y queremos mantener un equilibrio
entre el tamaño del conjunto de entrenamiento y el de test.
Para esta situación, se propone la siguiente variación 
del $t$-test. Se considera $d_i^{(j)}$, para $j=1,2$, $i=
1,\dots,5$, la diferencia entre el rendimiento de los dos
clasificadores para la repetición $i$ y el \textit{fold} $j$.
Se considera la media para cada repetición $\bar{d}_i = 
(d_i^{(1)} + d_i^{(2)})/2$ y la varianza estimada
$s_i^2 = (d_i^{(1)}-\bar{d_i})^2 + (d_i^{(2)}-\bar{d_i})^2$. 
Bajo la hipótesis nula de que los dos clasificadores tienen 
el mismo rendimiento, $d_i^{(j)}$ puede considerarse 
aproximadamente normal . Si suponemos $d_i^{(1)}$ y 
$d_i^{(2)}$ independientes y normales (lo cual no es
cierto porque los conjuntos de test y de entrenamiento
no se han obtenido de manera independiente), entonces
$s_i^2/\sigma^2$ sigue una distribución $\chi^2$ con un
grado de libertad. Si suponemos las $s_i^2$ independientes,
su suma $M = \frac{\sum\limits_{i=1}^5 s_i^2}{\sigma^2}
\sim \chi^2_5$ y 
\begin{equation}
	\label{eq:5x2CVt}
	t = \frac{d_1^{(1)}/\sigma}{\sqrt{M/5}} = 
		\frac{d_1^{(1)}}
			{\sqrt{\sum\limits_{i=1}^5 s_i^2/5}}
		\sim t_5.
\end{equation}


\paragraph{Comparación de algoritmos de clasificación} 
\cite{dietterich1998approximate}
En este artículo se estudia la comparación de algoritmos
aplicados a un único problema cuando disponemos de pocos
datos. Para el diseño y la evaluación de test estadísticos
se identifican las principales fuentes de 
variabilidad de los experimentos, como son los 
datos utilizados, la posible aleatoriedad de los 
propios algoritmos de aprendizaje y los posibles errores en
los datos disponibles. Se proponen los test de McNemar, un
test basado en él, usando el estadístico
	\[
		z = \frac{c_{01}-c_{10}}{\sqrt{2p(1-p)}}
	\]
	con $p$ el error medio de los dos clasificadores $\left(
p = \frac{2c_{00}+c_{01}+c_{10}}{2N}\right)$, que sigue
aproximadamente una distribución normal. Además,
se proponen las modificaciones del $t$-test con repetidas
muestras apareadas, el $t$-test pareado con validación
cruzada sobre$k$-\textit{fold} o el 5$\times$2 CV $t$-test.\\
	El experimento realizado consiste en, fijado un error 
$\varepsilon$, se le asigna a la mitad de los puntos
de un conjunto una probabilidad $\frac{\varepsilon}{2}$
de ser erróneamente clasificados
y a la otra mitad de los puntos la probabilidad
$\frac{3\varepsilon}{2}$ para uno de los algoritmos a
comparar. Para el otro algoritmo asignamos a la primera
mitad la probabilidad $\frac{3\varepsilon}{2}$ y a la
segunda mitad la probabilidad de fallar 
$\frac{\varepsilon}{2}$. Por tanto no deberíamos
rechazar la hipótesis nula, con lo que se realizan estos
experimentos para calcular el error de tipo I
para cada test. Se observa con la utilización de estos
datos generados sintéticamente que el error de tipo I más
elevado es para el $t$-test con repetidas
muestras apareadas.


\paragraph{Modificación sobre el 5$\times$2 CV $t$-test} Nótese que en la 
ecuación~\ref{eq:5x2CVt} el valor $d_1^{(1)}$ podría ser
ocupado por cualquier estadístico $d_i^{(j)}$. Alpaydin 
(\cite{Alpaydin98combined5})
propuso combinar estos valores de la forma:
	\[
		N = \frac{\sum\limits_{i=1}^5
				\sum\limits_{j=1}^2
					\left( d_i^{(j)} \right)^2}
				{\sigma^2}
			\sim \chi_{10}^2.
	\]
	Sustituyendo esto en \ref{eq:5x2CVt}, tenemos
un estadístico que es el cociente de dos v.a. que siguen 
una $\chi^2$, con lo que definimos la v.a. $f$, que seguirá
una distribución $F$ de Snedecor con 10 y 5 grados de
libertad:
	\[ 
		f = \frac{N/10}{M/5} \sim F_{10,5}.
	\]	
	
\section{Test no paramétricos}
 
\paragraph{Validación de un modelo de degradación de 
documentos} \cite{DBLP:journals/pami/KanungoHBSM00} En este artículo se presenta un 
modelo de la degradación en documentos producida por la 
digitalización de documentos impresos o la realización de 
fotocopias. El problema estadístico que se define es el 
siguiente. Dadas las muestras $x_1, \dots, x_N$ e $y_1, 
\dots, y_M$ de caracteres degradados de un mismo carácter de 
un documento y generados de manera artificial, 
respectivamente, realizar el test sobre la hipótesis nula de 
que ambos  provienen de la misma población para un valor de 
significación $\alpha$. Se aplica el algoritmo 
(\ref{alg:CMC-digitalizacion}) para realizar el test basado 
en permutaciones calculando el $p$-valor utilizando el 
algoritmo~\ref{alg:CMC-pvalue}.

\begin{algorithm}
	\caption{Test basado en permutaciones para validación en 
			modelo de degradación}
	\label{alg:CMC-digitalizacion}
	\begin{algorithmic}[1]
	\REQUIRE
		\begin{enumerate}[a]
		\item Datos reales $X = \{ x_1, \dots, x_N \}$
		\item Datos generados $Y = \{ y_1, \dots, y_M \}$
		\item Función distancia entre conjuntos $\rho(X,Y)$
		\item Función distancia entre caracteres $\delta(x,y)$
		\item Tamaño máximo del test $\alpha$
		\end{enumerate}
		\STATE Calcular $d_0 = \rho(X,Y)$
		\STATE Crear una muestra 
			$Z= \{x_1, \dots, x_N, y_1, \dots, y_M \}$.
		\FOR { $i \in \{1, \dots, K\}$ }
			\STATE Realizar una permutación de $Z$.
			\STATE Particionar $Z$ en $X'$ e $Y'$ tales que 
				$X' = \{ z_1, \dots, z_N \}$, $Y' = 
				\{ z_{N+1}, \dots, z_{N+M} \} $.
			\STATE Calcular $d_i = \rho(X,Y)$
		\ENDFOR
		\STATE Calcular la distribución empírica de las $d_i$
		\STATE Calcular el $p$-valor: 
			$\alpha_0 = P(d \geq d_0)$
		\STATE Rechazar la hipótesis nula si 
			$\alpha_0 \leq \alpha$
	\end{algorithmic}
\end{algorithm}

	
\paragraph{Comparación estadística de varios clasificadores} 
\cite{DBLP:conf/mlmta/ChenC03} En este artículo se describe un método 
para realizar la comparación estadística de varios 
clasificadores sobre una misma base de datos a través del 
test de Cochran, una generalización del test de McNemar. La 
aplicación de un test no paramétrico en lugar de uno 
paramétrico como ANOVA se justifica en el artículo debido a 
que no se puede asegurar la independencia entre las 
instancias clasificadas por un mismo clasificador, con lo que 
queda en entredicho el desempeño del test ANOVA, por ejemplo. 
Si la hipótesis nula, $H_0: \theta_1 = \dots = \theta_n$ 
(todos los algoritmos tienen un rendimiento equivalente), es 
cierta, no es necesario ningún análisis posterior. Si es 
rechazada, se pretenden realizar múltiples comparaciones para 
encontrar los algoritmos con mejor rendimiento. Para ello se 
describen los métodos de Bonferroni y Scheffé. El primero 
parte de una familia de $g$ contrastes, mientras que el 
segundo considera la familia $L$ de todos los contrastes 
posibles. Por ello, para el primero se deben prefijar las 
comparaciones a realizar, mientras que para el segundo esto 
se puede realizar \textit{a posteriori} según los datos 
obtenidos. Para escoger entre un método u otro, se recomienda 
escoger el primero cuando $Z_{1-\alpha/2g}$ (percentil $(1-
\alpha/2g)100$ de una distribución normal estándar) es menor 
que $\sqrt{\chi^2_{m-1,1-\alpha}}$ (donde $m$ es el número de 
algoritmos en la comparación) debido a que en este caso el 
método de Bonferroni tiene una mayor potencia. Escogeremos el 
método de Scheffé en caso contrario.

\paragraph{Comparación estadística de clasificadores en 
varios conjuntos de datos} \cite{DBLP:journals/jmlr/Demsar06} De la comparación 
utilizando un único conjunto de datos se pasa en este 
artículo a presentar un método para comparar el rendimento 
sobre varios conjuntos de datos utilizando test no 
paramétricos.\\
	Se incluyen en una primera sección las opciones para 
realizar las comparaciones entre dos clasificadores para 
varias conjuntos de datos, evaluando sus ventajas e 
inconvenientes:
	\begin{enumerate}
	\item Media entre los conjuntos de datos: Arroja poca 
		información poco valiosa al combinar puntuaciones 
		obtenidas en distintos dominios.
	\item $t$-test emparejados: El estadístico de este test 
		paramétrico es similar al usado al realizar la media 
		entre los conjuntos de datos. Además, se necesita que 
		la diferencia siga una distribución normal, además de 
		que los datos disponibles suelen ser reducidos. 
	\item Test de signo: Como se ha dicho en el desarrollo de 
		los test, no tiene en cuenta la magnitud de la 
		diferencia.
	\item Test de rangos con signo de Wilcoxon: Más seguro 
		que el $t$-test al no suponer la distribución normal. 
		Si no se dan las circunstancias adecuadas, este test 
		puede tener mayor potencia que el $t$-test.
	\end{enumerate}
	
	En cuanto a la comparación entre varios clasificadores, 
se incluye una breve descripción del método ANOVA y una 
descripción más extensa del test de Friedman. Se incluye el 
estadístico $F_F = \frac{(n-1) \chi^2_F}{n(k-1) - \chi^2_F}$ 
que sigue la distribución $F$ con $k-1$ y $(k-1)(n-1)$ grados 
de libertad. Aunque el test ANOVA tiene mayor potencia cuando 
las condiciones se cumplen, se incluye un experimento que 
muestra que apenas hay diferencias si éstas no se dan.\\

	Una vez se ha rechazado la hipótesis nula se describen 
los test \textit{post-hoc} de Nemenyi, Bonferroni-Dunn, Holm, 
Hochberg y Hommel.
	
	
\paragraph{Extensión sobre los procedimientos 
\textit{post-hoc}} \cite{garcia2008extension} Este artículo 
extiende al anterior prestando una mayor atención a la 
descripción de los procedimientos para la comparación entre 
algoritmos una vez rechazada la hipótesis nula. Se introducen
los procedimientos de Nemenyi y Holm. Debido a que las 
posibles hipótesis sobre la comparación  por parejas están 
relacionadas y no se pueden dar todas las posibles 
combinaciones, se presentan los procedimientos de 
Schaffer estático y dinámico. Bergmann y Hommel propusieron 
un procedimiento basado en la idea de encontrar todas las 
hipótesis que no pudieran ser rechazadas. En un estudio 
comparativo sobre el rendimiento de cinco clasificadores 
sobre treinta conjuntos de datos se observa como el 
procedimiento de Bergmann y Hommel tiene una mayor 
potencia.\\

	Se aborda también el problema de calcular $p$-valores 
ajustados (APV) que tengan en cuenta los demás test 
realizados para poder utilizar directamente este valor 
ajustado. Se indican el cálculo para cada procedimiento 
\textit{post-hoc}. Tras realizar un nuevo estudio 
introduciendo el uso de APV, se descarta el test de Nemenyi 
debido a que es demasiado conservador. Se recomienda el 
procedimiento estático de Schaffer frente al del Holm al 
tener mayor potencia al usar las relaciones entre las 
hipótesis y una dificultad no muy superior. Debido a su coste 
computacional y su complejidad, se recomienda el método de 
Bergmann-Hommel sólo para cuando las diferencias entre 
clasificadores sean muy pequeñas y otros métodos no detecten 
diferencias significantivas.

	\paragraph{Estudio del uso de test no paramétricos} 
	\cite{DBLP:journals/heuristics/GarciaMLH09} En este artículo se realiza un estudio
sobre el uso de test no paramétricos para comparar el 
rendimiento de algoritmos evolutivos en la Sesión Especial de 
Optimización de parámetro real en el Congreso de Computación 
Evolutiva (CEC 2005). Los datos disponibles son los errores 
cometidos por los algoritmos a la hora de minimizar 25 
funciones de 10, 30 y 50 variables reales. En primer lugar se 
estudian las condiciones para realizar un test paramétrico 
($t$-test) llevando a cabo test de normalidad (test de 
Kolmogoro-Smirnov, test de Shapiro-Wilk y el de D'Agostino-
Pearson) y el test de Levene de homocedasticidad. Se 
pretende realizar en primer lugar la comparación entre dos de 
los algoritmos para comprobar si hay diferencias entre ellos. 
Sin embargo los test de normalidad rechazan para los 
resultados en casi todas las funciones a minimizar la 
hipótesis de normalidad y el test de Levene rechaza la 
hipótesis de la equivalencia de las varianzas, con lo que no 
se cumplen las condiciones para realizar el $t$-test. Si se 
realizan el $t$-test y el test no paramétrico de Wilcoxon se 
observa que los resultados son muy similares.\\
	A la hora de afrontar la comparación en múltiples 
problemas en lugar de problema por problema en este artículo 
se toma la media de las distintas ejecuciones. Al realizar 
los test de normalidad sobre las dos muestras de 25 elementos 
para cada algoritmo se descarta la normalidad. El $t$-test y 
el test de Wilcoxon no rechazan la hipótesis nula de que 
ambos algoritmos sean igual de buenos, sin embargo el $p$-
valor del test de Wilcoxon es bastante menor, lo que nos 
lleva a pensar que, debido a la ausencia de normalidad y de 
pocos datos y difíciles de incrementar (pues en este caso 
deberíamos añadir nuevas funciones y sus 
correspondientes evaluaciones para cada algoritmo), en estos 
casos es mejor utilizar test no paramétricos.\\
	Para realizar el análisis sobre todos los algoritmos se 
separan las funciones fáciles, aquellas en las que algún 
algoritmo alcanza el óptimo frente a las que ningún algoritmo 
llegó al óptimo. El test no paramétrico de Friedman y el de 
Iman-Davenport arrojan valores que nos llevan a rechazar la 
hipótesis nula tanto usando el grupo de funciones difíciles 
como usando el conjunto total de funciones.\\
	El estudio posterior se centra en comparar los algoritmos 
con aquel que consiguió un menor error medio, el 
\textit{G-CMA-ES}. Para ello se utilizan el test de 
Bonferroni-Dunn, el de Holm y el de Hochberg para los dos 
grupos de funciones. El test de Hochberg llega a encontrar 
diferencias entre el mejor algoritmo y todos los demás para 
$\alpha=0.1$. 
	
	\paragraph{Estudio estadístico sobre algoritmos 
genéticos: precisión e interpretabilidad} \cite{Garcia2008} En 
este artículo se realiza la comparación sobre algoritmos 
evolutivos utilizados en clasificación, comparando tanto la 
tasa de clasificación como el valor kappa de Cohen y la 
interpretabilidad de estos resultados. Al igual que en el 
artículo anterior, se realizan test de normalidad y 
homocedasticidad para comprobar si se dan las condiciones 
para aplicar un test paramétrico. Si las comparaciones por 
parejas de problemas se repiten crece el error asociado FWER 
(\textit{family-wise error rate}), esto es la probabilidad de 
cometer al menos un error en la familia de hipótesis. Para realizar 
comparaciones entre varios algoritmos y fijar el FWER de 
antemano se utilizan el test de Friedman y el de Iman-
Davenport y los test de Bonferroni-Dunn y el de Holm para 
buscar diferencias entre los algoritmos una vez descartada la 
hipótesis de que todos tienen la misma media. Se 
calculan también los $p$-valores ajustados.\\
	Para realizar el análisis de la interpretabilidad de los 
resultados se considera $\textit{Complejidad} = 
\textit{Tamaño} \cdot ANT$. Se realiza para este estudio 
únicamente un procedimiento para múltiples comparaciones, 
obteniendo diferencias significativas con el test de 
Bonferroni-Dunn y nivel de significación $\alpha=0.1$.
	
	\paragraph{Test avanzados para la comparación y análisis 
experimental de la potencia} \cite{DBLP:journals/isci/GarciaFLH10} En este
artículo se introducen con respecto a otros artículos 
previos los test de múltiples signos y la estimación del 
contraste basado en medianas como test básicos para realizar
entre múltiples clasificadores y conjuntos de datos como 
test básicos, y los más avanzados del test de Friedman
de rangos alineados (\textit{Friedman Aligned Ranks}, 
Friedman-AR) y el test de Quade. Además, para realizar las 
comparaciones  posteriores  una vez descartada la hipótesis 
de un igual  rendimiento entre todos los algoritmos, se 
incluyen a los  procedimientos \textit{post-hoc} los de 
Holland y Finner  dentro de  los métodos descendentes, el de 
Rom como método  ascendente y el procedimiento de Li, para 
rechazar hipótesis en dos pasos.\\
	Se realiza un estudio sobre la potencia de los test
de Friedman, de Friedman-AR y de Quade, utilizando 8 y 4
algoritmos sobre 10 y 6 conjuntos de datos respectivamente. 
Se observa que el test de Friedman-AR es más conservador
cuando el número de algoritmos es elevado. En este estudio
se observa que el test de Quade se comporta mejor, sin
embargo depende de los conjuntos de datos utilizados. Para
el estudio sobre los test \textit{post-hoc}, se utilizan
8 algoritmos sobre 10 conjuntos de datos seleccionados
de forma pseudo-aleatoria (regulando mediante un parámetro la 
probabilidad de seleccionar una base de datos según la 
diferencia de los algoritmos a comparar en ella). Los 
resultados obtenidos hacen desaconsejar el test de 
Bonferroni-Dunn, en un siguiente grupo estarían Holm, 
Hochberg, Hommel, Holland y Rom con una potencia similar
y los que mejores resultados obtienen son el test de Finner
y el test de Li, el cual ofrece mejores resultados 
cuando las diferencias entre el algoritmo de control y los 
demás sean grandes.

	
	\paragraph{Análisis \textit{bootstrap} de múltiples 
repeticiones de experimentos} \cite{Otero201488} En este artículo 
se presenta la realización de test basados en permutaciones 
para la comparación del rendimiento de algoritmos de 
clasificación, realizando un estudio para el que el test 
presentado tiene mayor potencia que los test de ANOVA o 
Friedman. El contexto para el que se presenta este método es 
el de la comparación entre algoritmos estocásticos (de ahí 
que se realice la repetición de la ejecución de los 
algoritmos con diferentes semillas) utilizando validación 
cruzada. Se propone usar intervalos para recoger parte de la 
variabilidad de los datos atribuible a la repetición del 
experimento. Se plantea el problema del uso de la media entre 
los distintos experimentos para evaluar un algoritmo debido a 
que los valores extremos tienen una influencia excesiva y se 
aboga por otros estadísticos como la mediana, el rango 
intercuartílico o la media restringida a los valores 
centrales.\\
	El primer test es un test basado en permutaciones. Parte 
de la muestra compuesta por $e_{adfr}$, que es el error del 
algoritmo $a$-ésimo, el conjunto de datos $d$, para el 
\textit{fold} $f$ y la repetición $r$. Se utiliza en este 
test como estadístico para evaluar cada algoritmo la media de 
las ejecuciones en todos los \textit{folds}, repeticiones y 
conjuntos de datos como estimador, utilizando el algoritmo ~
\ref{alg:CMC-pvalue}. Se rechazará la hipótesis nula (la 
equivalencia de las medias) si para algún $a$ la media 
inicial $\hat{e}_a$ se encuentra en las colas de la 
distribución formada por los valores de las permutaciones $
\hat{e}_a^*$ consistente en la cdf muestral:
	\[ 
		\hat{F}_{a \cdot \cdot \cdot}(x) =
			\frac{1}{n_d n_r n_f}
			\sum\limits_{i=1}^{n_d}
				\sum\limits_{j=1}^{n_f}
					\sum\limits_{k=1}^{n_r}
						\mathbb{I}[ e_{aijk} \leq x ].
	\]
	El segundo test sí hace uso de los intervalos, notando 
por $[q_{-adf}, q_{+adf}]$ el intervalo de los errores para 
un algoritmo, base de datos y \textit{fold}. Se definen las 
cdf muestrales como
	
\begin{align*}
	\hat{F}_{-a \cdot \cdot}(x) &=
		\frac{1}{n_d n_f}
			\sum\limits_{i=1}^{n_d}
				\sum\limits_{j=1}^{n_f}
					\mathbb{I}[ q_{+aij} \leq x ], \\
	\hat{F}_{+a \cdot \cdot}(x) &=
		\frac{1}{n_d n_f}
			\sum\limits_{i=1}^{n_d}
				\sum\limits_{j=1}^{n_f}
					\mathbb{I}[ x \in [q_{-aij},q_{+aij}) ] +
			\hat{F}_{-a \cdot \cdot}(x) .			
\end{align*}

	En este test se aplica también el algoritmo 
~\ref{alg:CMC-pvalue}, rechazando la hipótesis nula si el 
intervalo $[\hat{q}_{-a}, \hat{q}_{+a}]$, con $\hat{q}_{-a}, 
\hat{q}_{+a}$ las medias de los extremos inferiores y 
superiores respectivamente para cada algoritmo con respecto a 
las conjuntos de datos y \textit{fold}, se encuentra en las colas 
de la distribución de los $[\hat{q}_{-a}^*, \hat{q}_{+a}^*]$.
\\
	Los resultados incluidos ofrecen una comparación entre la 
potencia de estos test presentados, el test de Friedman y el 
de ANOVA. Se realiza un experimento con datos sintéticos 
consistentes en los resultados de los errores de 5 algoritmos 
en 32 conjuntos de datos para una secuencia de diferencias en sus 
medias preestablecidas. Para calcular la potencia se estima 
para cada diferencia el porcentaje de experimentos para el 
que se encuentran diferencias significativas entre los 
algoritmos dos a dos. Los resultados indican una mayor 
potencia de los test presentados en este artículo.


\section{Test bayesianos}

	Se presenta en este apartado una revisión del estado del 
arte con diferentes métodos para la comparación de algoritmos
de aprendizaje automático mediante la realización de test 
estadísticos. 

\paragraph{Versión bayesiana del test de Wilcoxon} 
\cite{DBLP:conf/icml/BenavoliCMZR14} Para la comparación de algoritmos mediante
muestras apareadas, se presenta el test de Wilcoxon en 
versión bayesiana. Se realiza, al igual que en el test
de signo y el de rangos con signo, utilizando el Proceso de
Dirichlet para la distribución \textit{a priori}. La versión 
bayesiana, a diferencia de la no paramétrica, no necesita 
suponer la simetría de las muestras. Una cuestión importante
es la elección de $G_0$ en caso de que no dispongamos
información. El enfoque realizado se basa en la obtención
del límite de $DP$ para $s \rightarrow 0$. Otro enfoque 
consiste en modelar la falta de información \textit{a priori} 
como un conjunto de distribuciones de probabilidad. 

\paragraph{Procedimiento no paramétrico bayesiano para la 
comparación de algoritmos} \cite{DBLP:conf/icml/BenavoliCMZ15} En este 
artículo se presenta la versión bayesiana del test de 
Friedman. Se presenta el test, que utiliza el $DP$ para 
realizar la comparación de dos algoritmos en un primer
momento y la versión para tres o más. Una vez se ha rechazado 
la hipótesis de que los algoritmos se comporten de igual
manera, se procede realizando la comparación dos a dos 
para encontrar dónde se producen las diferencias. Se realiza 
también en este artículo un experimento comparando el test
bayesiano de Friedman con el test $F$-\textit{race}, test que
compara varios algoritmos según aquel que consiga el mejor
resultado.

\paragraph{Test estadístico para el análisis conjunto de 
medidas de rendimiento} \cite{DBLP:conf/jsai/BenavoliC15}	Para la 
comparación de algoritmos con respecto a varias medidas
se suelen utilizar o bien la media ponderada de las medidas
para cada algoritmo o bien el frente de Pareto, donde se 
consideran las soluciones no dominadas (aquellas que no son
peor con respecto a otra solución para todos los criterios).
En este artículo se presenta en primer lugar la realización 
de un test basado en la razón de verosimilitud, que presenta
las mismas desventajas que los THN. Para la realización del
test bayesiano, se opta por usar como distribución 
\textit{a priori} la distribución de Dirichlet, obteniéndose
también como distribución \textit{a posteriori} la 
distribución de Dirichlet de la que se calculan las 
probabilidades deseadas mediante muestras aleatorias.
Finalmente, se presenta una mejora de este método usando
redes bayesianas.

\paragraph{Comparación de algoritmos en múltiples bases de 
datos con CV} \cite{DBLP:journals/ml/CoraniB15} En este artículo se introduce
el $t$-test bayesiano (\ref{ssec:bayes-ttest}) y un test de 
Poisson que compara dos clasificadores en varias bases de 
datos teniendo en cuenta la correlación y la incertidumbre 
provocada por la validación cruzada en cada conjunto de 
datos.

\paragraph{Tutorial para la comparación de clasificadores 
mediante análisis bayesiano}. \cite{DBLP:journals/corr/BenavoliCDZ16} Este artículo 
es un artículo divulgativo que presenta las ventajas de los 
test bayesianos y expone su uso en diferentes casos, 
aplicando el $t$-test bayesiano, el test de signos y el $t$-
test bayesiano jerárquico (Corani, \cite{coranistatistical}), una extensión que permite realizar 
inferencia en varios conjuntos de datos.

\chapter{Software}
\label{chapter:software}

\section{Revisión software existente}

	En esta sección se incluye una descripción del software 
integrado en R.

\subsection*{JavaNPST}
	 
	Este paquete desarrollado por Derrac \textit{et. al} 
(\cite{2015arXiv150104222D}) implementado en \texttt{Java} 
contiene cuarenta test estadísticos no paramétricos de las 
siguientes familias:
\begin{itemize}
	\item Test de aleatoriedad
	\item Test de bondad del ajuste
	\item Procedimientos de una muestra y de muestras
		 apareadas
	\item Procedimientos generales de dos muestras
	\item Test para problemas de localización
	\item Test para problemas de escala
	\item Test de igualdad de muestras independientes
	\item Medidas de asociación para muestras bivariantes
	\item Medidas de asociación en clasificaciones múltiples
	\item Análisis del conteo de datos
\end{itemize}
	
	La biblioteca \texttt{JavaNPST} incluye una definición de 
las estructuras de datos (tablas y secuencias) usadas para 
el manejo y la realización de los test incluidos. En el 
artículo se incluyen un ejemplo de uso aplicado a la 
comparación de algoritmos mediante los test de 
signo y el test de Wilcoxon. 

\subsection*{scmamp} 

	Este paquete de \texttt{R} desarrollado por Calvo y 
Santafe (\cite{scmamp}) implementa test estadísticos para
la comparación de múltiples algoritmos en distintos 
problemas. Una de los principales atractivos radica en la 
disponibilidad de ejecutar los test \textit{post hoc} que
constituyen el estado del arte. Podemos encontrar los 
siguientes tipos de funciones en este paquete:

\begin{description}
	\item[Test para la comparación de múltiples algoritmos:]
		Disponemos del test paramétrico ANOVA y de test
		no paramétricos como el contraste basado en
		medianas, el test de Friedman, el de rangos alineados
		de Friedman o el test de Quade. Cada uno de ellos
		tiene un método adicional para efectuar el test
		 \textit{post hoc} asociado.
	\item[Ajuste de $p$-valores] Se incluyen los métodos
		para ajustar los $p$-valores descritos por García en
		\cite{garcia2008extension} y 
		\cite{DBLP:journals/isci/GarciaFLH10}:
		Bergmann-Hommel, Finner, Holland, Li, Rom y Shaffer.
		Se incluye una función para obtener los $p$-valores
		de las comparaciones que deseemos realizar entre
		algoritmos.
	\item[Conjuntos de datos] utilizados en distintos
		artículos sobre los que se realizan los experimentos. 
	\item[Métodos auxiliares y creación de gráficos] La
		representación de los resultados se facilita 
		con las funciones incluidas en el paquete.
\end{description}

	Junto con el paquete viene un tutorial para saber cómo
se manejan las funciones y cómo utilizar el paquete para
la comparación del rendimiento de algoritmos.

\subsection*{BayesianML-tutorial}

	Por último, presentamos los test bayesianos implementados
por Benavoli (\cite{DBLP:journals/corr/BenavoliCDZ16}) y 
disponibles en el repositorio \fnurl{\texttt{tutorial}}
{https://github.com/BayesianTestsML/tutorial/}. 
En este repositorio se encuentra la implementación en 
\texttt{R} para las versiones bayesianas de los test 
estadísticos presentados en este artículo, es decir, 
el test de Wilcoxon de rangos con signo, el test de signo
y el $t$-test correlado bayesiano. Se incluye también el
test jerárquico para la comparación sobre varios conjuntos
de datos descrito por Corani \textit{et. al} en \cite{coranistatistical}. 
El código para este método se ha realizado
en \texttt{R} y \texttt{Stan}, un lenguaje para inferencia
bayesiana. Para ilustrar la comparación, en el repositorio
se encuentran disponibles los tests en el lenguaje \texttt{Julia} 
y una serie de \textit{Notebooks} 
de este lenguaje para ilustrar su uso, acompañándolo
con gráficas de la distribución \textit{a posteriori}
de la media de las diferencias.
	
	 
\section{rNPBST}

	En esta sección se describirá la herramienta software
desarrollada. El objetivo inicial planteado consistía en
la integración en un paquete de los test estadísticos
más utilizados para la comparación de algoritmos. Para
esta labor, partimos de los paquetes descritos.

\subsection*{JavaNPST}

	La integración del paquete \texttt{JavaNPST} en
\texttt{R} se ha realizado usando \texttt{rJava}.\\
Para ello, en primer lugar se definen en \texttt{R} los 
métodos necesarios para pasar de los tipos de datos 
habituales como las matrices y los vectores a las estructuras 
usadas en el paquete  \texttt{JavaNPST}: \texttt{dataTable}, 
que son el equivalente a las matrices, y la clase abstracta 
\texttt{Sequence}  de la que heredan las clases 
\texttt{NumericSequence} y \texttt{StringSequence}, que se 
corresponden con vectores, bien de números o bien de cadenas 
de caracteres.\\
	Los test disponibles en el paquete \texttt{JavaNPST} 
se corresponden con una clase de la cual se crea un 
objeto y se le proporcionan unos datos, entonces realizamos
el test y el objeto creado nos devuelve una serie de valores.
Estos valores se corresponden con los estadísticos utilizados
en la realización del test o los $p$-valores devueltos.
La forma de llevar esto a \texttt{R} ha sido mediante
una función para cada test, donde se crea un objeto
de la clase correspondiente con los datos (bien una matriz
o una secuencia), se ejecuta el test, obteniendo el informe
generado por el paquete utilizado y finalmente se obtienen 
los estadísticos utilizados en el test y los $p$-valores
disponibles (según el test tendremos un $p$-valor o 
varios según las hipótesis alternativas). Estos dos últimos
datos, los estadísticos y los $p$-valores los introducimos
en un objeto de la clase de \texttt{R} \texttt{htest},
que frecuentemente se usa para guardar los resultados de un
test estadístico.

\subsection*{BayesianML-tutorial}

	Como se ha descrito, en este repositorio están
implementados los test bayesianos en \texttt{R}
y en \texttt{Julia}, sin embargo, la información
que devuelven es distinta: en la versión en \texttt{R}
únicamente se devuelve la probabilidad \textit{a
posteriori} de la pertenencia del parámetro 
$\theta$ la diferencia entre los clasificadores
a la región izquierda, a la \textit{rope} o a la
región derecha; mientras que en la versión en
\texttt{Julia} se devuelven una muestra de 
la distribución de Dirichlet \textit{a posteriori},
que facilita la representación gráfica.\\
Por tanto, se ha descartado la opción de 
llamar a estas funciones desde nuestro paquete
(si quisiéramos una representación gráfica
y la probabilidad de pertenencia estaríamos
realizando dos veces el test) y optamos por
escribir nuestra propia función para el
test bayesiano de signo, el test bayesiano de 
rangos con signos y la versión bayesiana 
de $t$-test correlados. Además, se ha observado 						%TODO Revisar
un tiempo de ejecución excesivo utilizando
el test bayesiano de rangos con signos (en torno
a dos minutos) en la versión del repositorio
frente a la versión implementada (unos cinco
segundos). Sin entrar en la implementación, se 
describe a continuación el algoritmo utilizado:

\begin{algorithm}[H]
	\caption{Test bayesiano de rangos con signo}
	\label{alg:TBRS}
	\begin{algorithmic}[1]
	\REQUIRE $x, y$, vectores de observaciones.
	\REQUIRE $s, z_0, weights$, parámetros \textit{a priori}.
	\REQUIRE $rope_{min}, rope_{max}$, extremos de la \textit{rope}
	\REQUIRE $n_{samples}$ Número de muestras a obtener de
		la distribución de los parámetros.
	\STATE $d \leftarrow$ Diferencia entre los vectores 
		de observaciones 
	\STATE $pesos \leftarrow$ muestra aleatoria de 
		la distribución de Dirichlet de tamaño $n_{samples}$.
	\STATE $b_l, b_{rope}, b_r \leftarrow$ vector binario con
		la pertenencia o no de cada elemento de $d$ a la 
		región izquierda, a la \textit{rope} y a la derecha.
	\FOR {$i \in [1, n_{samples}]$}
		\STATE $coef \leftarrow pesos_i^T pesos_i$
		\STATE $muestra_i \leftarrow$
			$(coef \cdot b_l, coef \cdot b_{rope}, 
			coef \cdot b_r)$,  
	\ENDFOR
	\STATE $probabilidades_{a\ posteriori} \leftarrow $
		Media del número de veces que el máximo de la fila 
		está en cada columna.
	\RETURN $muestra, probabilidades_{a\ posteriori}$
	\end{algorithmic}
\end{algorithm} 


Para la representación gráfica se ha utilizado
la representación que vemos en el artículo
de Benavoli (\cite{DBLP:journals/corr/BenavoliCDZ16}),
\texttt{R} utilizando el paquete \texttt{ggtern} (\cite{ggtern})
para la representación de ternas. Estos gráficos
muestran la nube de puntos obtenida de realizar
una muestra aleatoria de una distribución
de Dirichlet con los parámetros obtenidos.
El centro del triángulo representa el punto
$\left( \frac{1}{3}, \frac{1}{3}, \frac{1}{3} \right)$
es decir, hay una igual probabilidad para que el punto
caiga en la \textit{rope}, de un lado u otro. 



\chapter{Experimentos}
\label{chapter:experimentos}
\section{Introducción al caso práctico}
	Una vez hemos presentado las principales técnicas
estadísticas que se utilizan en la comparación de algoritmos
en aprendizaje automático, realizaremos un estudio práctico
en \texttt{R} sobre la aplicación de estas técnicas.\\
	Para ello, usaremos los resultados obtenidos por cinco
algoritmos (Tabla~\ref{tab:algoritmos}) 
sobre varios conjuntos de datos (Tabla~\ref{tab:DS}). Los
conjuntos de datos constituyen un subconjunto de los 
disponibles en el repositorio 
\fnurl{KEEL}{http://www.keel.es/} 
(\cite{alcala2010keel}) para problemas de clasificación
supervisada. En dicho repositorio encontramos ya disponibles
particiones de los conjuntos de datos para facilitar que los
experimentos sean reproducibles. Se ha utilizado para cada 
conjunto de datos la partición para realizar 10-CV, esto es,
cada conjunto se particiona en 10 subconjuntos y en cada una 
de las 10 repeticiones, utilizamos un subconjunto como test y
los restantes 9 como datos de entrenamiento. Cada archivo
contiene los diez pares de ficheros con datos de
entrenamiento y test. Los algoritmos a comparar son
ampliamente conocidos y no se entra en más detalle sobre
ellos ya que la materia de este estudio no son los algoritmos
sino los test utilizados para compararlos.

\begin{table}[H]
\centering
\caption{Algoritmos}
\label{tab:algoritmos}
\begin{tabularx}{\textwidth}{lX}
\toprule
Algoritmo             & Descripción                                                                                                                             \\ \midrule
\texttt{multinom}     & Regresión logística de la biblioteca \texttt{nnet} para un problema de clasificación con varias clases.                               \\
\texttt{knn}          & Algoritmo de los $k$ vecinos más cercanos de la biblioteca \texttt{class}. 
Se han utilizado los valores por defecto $k=1$, $l=0$ y la distancia euclídea.      \\
\texttt{randomForest} & Algoritmo de la biblioteca \texttt{randomForest}. Se ha fijado el parámetro \texttt{mtray} al valor por defecto $mtray = \sqrt{p}$. \\
\texttt{nnet}         & Red neuronal de la biblioteca \texttt{nnet}.                                                                                          \\
\texttt{naiveBayes}   & Clasificador \textit{naive Bayes} de la biblioteca \texttt{e1071}.                                                                  \\ \bottomrule
\end{tabularx}
\end{table}

\begin{table}[H]
\centering
\caption{Conjuntos de datos}
\label{tab:DS}
\begin{tabular}{@{}lccc@{}}
\toprule
Conjunto de datos & N$\degree$ de atributos & N$\degree$ de clases & N$\degree$ de instancias \\ \midrule
abalone           & 8                   & 28               & 4174                 \\
australian        & 14                  & 2                & 690                  \\
automobile        & 25                  & 6                & 205                  \\
balance           & 4                   & 3                & 625                  \\
breast            & 9                   & 2                & 277                  \\
bupa              & 6                   & 2                & 345                  \\
car               & 6                   & 4                & 1728                 \\
cleveland         & 13                  & 5                & 303                  \\
crx               & 15                  & 2                & 690                  \\
dermatology       & 34                  & 6                & 366                  \\
german            & 20                  & 2                & 1000                 \\
glass             & 9                   & 7                & 214                  \\
hayes-roth        & 4                   & 3                & 160                  \\
heart             & 13                  & 2                & 270                  \\
ionosphere        & 33                  & 2                & 351                  \\
led7digit         & 7                   & 10               & 500                  \\
letter            & 16                  & 26               & 20000                \\
lymphography      & 18                  & 4                & 148                  \\
mushroom          & 22                  & 2                & 8124                 \\
optdigits         & 64                  & 10               & 5620                 \\
satimage          & 36                  & 7                & 6435                 \\
spambase          & 57                  & 2                & 4597                 \\
splice            & 60                  & 3                & 3190                 \\
tic-tac-toe       & 9                   & 2                & 958                  \\
vehicle           & 18                  & 4                & 846                  \\
vowel             & 13                  & 11               & 990                  \\
wine              & 13                  & 3                & 178                  \\
yeast             & 8                   & 10               & 1484                 \\
zoo               & 16                  & 7                & 101                  \\ \bottomrule
\end{tabular}
\end{table}

\section{Aplicación de test estadísticos}

\section{Comparación de test estadísticos}
%
%%\documentclass[11pt,leqno]{book}
%\usepackage[spanish,activeacute]{babel}
%\usepackage[utf8]{inputenc}
%\usepackage{enumerate}
%
%\begin{document}


\chapter*{Conclusiones y trabajos futuros}\addcontentsline{toc}{chapter}{Conclusiones}


	Las conclusiones obtenidas con la realización de este trabajo son principalmente que la estadística es una herramienta fundamental para la comparación entre el rendimiento de algoritmos en el ámbito del aprendizaje automático y que la comparación de estos algoritmos debe 
realizarse de manera rigurosa para poder obtener resultados confiables. Se han observado cómo existen diversos métodos para la realización de los test estadísticos, los cuales 
deben usarse bajo unas ciertas condiciones que, en caso 
de no cumplirse, conllevan una pérdida de potencia en la comparación, como es el caso de los test paramétricos. Por otra parte, frente a los test de hipótesis nula, muchas veces malinterpretados, surgen los test bayesianos como solución a algunos de estos problemas.\\
Otra conclusión a destacar es la importancia del manejo de
los conjuntos de datos, la separación en conjuntos de entrenamiento y test y la validación cruzada.\\
Como trabajo futuro, queda pendiente la integración completa del paquete de Java y la profundización en nuevos test para la comparación de algoritmos. 


%\end{document}
%
%%\nocite{*}
%\begin{thebibliography}{99}

\bibitem{JAPSHA11}
  Nathalie Japkowicz, Mohak Shah,
  \emph{Evaluating Learning Algorithms - A Classification Perspective},
  Cambridge University Press,
  2011.

\bibitem{SANTAFE15}
  Guzmán Santafé, Iñaki Inza, José A. Lozano,
  \emph{Dealing with the evaluation of supervised classification algorithms},
  Springer Science - AIRE,
  2015.
  
\bibitem{PRATI11}
  Ronaldo C. Prati, Gustavo E.A.P.A. Batista, Maria Carolina Monard,
  \emph{A Survey on Graphical Methods for Classification Predictive Performance Evaluation},
  IEEE TRANSACTIONS ON KNOWLEDGE AND DATA ENGINEERING, VOL. 23, NO. 11,
  2011

\bibitem{DICCHA11}
  Jean Dickinson Gibbons, Subhabrata Chakraborti,
  \emph{Nonparametric Statistical Inference},
  Fifth edition,
  Chapman and Hall/CRC,
  2011
	
\bibitem{SHESKIN11}
  David J. Sheskin,
  \emph{Handbook of Parametric and Nonparametric Statistical Procedures},
  Fourth edition,
  Chapman and Hall/CRC,
  2007
  
\end{thebibliography}
%
%\appendix
%\input{apendices/manual_usuario/manual_usuario}
%%\input{apendices/paper/paper}
%\input{glosario/entradas_glosario}
% \addcontentsline{toc}{chapter}{Glosario}
% \printglossary
\chapter*{}
\thispagestyle{empty}
\bibliography{capitulos/bibliografia}\addcontentsline{toc}{chapter}{Bibliografía}
\label{bibliografia}
\bibliographystyle{plain}

\end{document}